\section{The projection chart}\label{sec:chart}

Absorbing the weights into the vectors replaces a design by a projection
on the index space, and replaces the domination condition
\eqref{eq:dominate} by a comparison between a principal block of that
projection and a diagonal matrix.

Throughout, $k\ge1$ and $m\ge k$; the second inequality is
Lemma~\ref{lem:span}.  For an index set $C$ of size $r$,
listed in increasing order, $M_C$ denotes the $r\times k$ matrix whose
rows are $g_c\adj$, $c\in C$, and
\[
  \Gram_C\;:=\;M_C^{\vphantom{\mathsf{H}}}M_C\adj
  \;=\;\bigl(\langle g_c,g_d\rangle\bigr)_{c,d\in C}
\]
their Gram matrix.  The two products of $M_C$ carry different
information:
\begin{equation}\label{eq:twoproducts}
  M_C\adj M_C^{\vphantom{\mathsf{H}}}\;=\;\sum_{c\in C}g_c^{\vphantom{\mathsf{H}}}g_c\adj
  \;=\;\SC \quad (k\times k),
  \qquad
  M_C^{\vphantom{\mathsf{H}}}M_C\adj\;=\;\Gram_C \quad (r\times r).
\end{equation}
Both identities are the definition of matrix multiplication read in the
two possible orders.

\subsection{Construction}\label{sec:chartconstruction}

Put
\[
  v_c\;:=\;\sqrt{t_c}\,g_c\in\FF^{k},
\]
and let $V$ be the $m\times k$ matrix whose $c$-th row is $v_c\adj$; thus
$V=\diag(\sqrt{t_1},\dots,\sqrt{t_m})\,M_{\{1,\dots,m\}}$.  The Parseval
condition \eqref{eq:parseval} is exactly the orthonormality of the
columns of $V$:
\begin{equation}\label{eq:isometry}
  V\adj V\;=\;\sum_{c}v_c^{\vphantom{\mathsf{H}}}v_c\adj
  \;=\;\sum_c t_c\,g_c^{\vphantom{\mathsf{H}}}g_c\adj\;=\;I_k .
\end{equation}
Accordingly
\begin{equation}\label{eq:chart}
  P\;:=\;VV\adj
\end{equation}
is a matrix on the \emph{index} space $\FF^m$.  Write
$D:=\diag(t_1,\dots,t_m)$ and
\begin{equation}\label{eq:chartgap}
  W\;:=\;P-D .
\end{equation}
The pair $(P,t)$ is the \emph{chart} of the design and $W$ its
\emph{gap}.

\begin{proposition}[Lean]\label{prop:chartid}
$P$ is the orthogonal projection onto the column space of $V$, of rank
$k$, and
\[
  P_{cd}=\sqrt{t_ct_d}\,\langle g_c,g_d\rangle,
  \qquad P_{cc}=s_c,\qquad \tr P=k .
\]
Consequently $W_{cc}=t_c(\ell_c-1)$ and $\sum_c W_{cc}=k-1$.
\end{proposition}

\begin{proof}
$P\adj=(VV\adj)\adj=VV\adj=P$, and, by \eqref{eq:isometry},
$P^2=V(V\adj V)V\adj=VIV\adj=P$.  A Hermitian idempotent is the orthogonal
projection onto its range; that range lies in the column space of $V$, and
$VV\adj V=V(V\adj V)=V$ puts the column space back inside it, a second use
of \eqref{eq:isometry}.  The rank of $P$ is therefore the rank of $V$,
which is $k$ because $V\adj V=I_k$.  The entries are
$P_{cd}=\langle v_c,v_d\rangle=\sqrt{t_ct_d}\langle g_c,g_d\rangle$; at
$c=d$ this is $t_c\ell_c=s_c$.  Finally
$\tr P=\tr(VV\adj)=\tr(V\adj V)=\tr I_k=k$.  The diagonal of $W$ is
$s_c-t_c=t_c(\ell_c-1)$, and $\sum_c(s_c-t_c)=k-1$ by
Lemma~\ref{lem:trace}.
\end{proof}

\begin{lemma}\label{lem:realize}
Fix $m$, $k$ and a weight vector $t$ with $t_c>0$ and $\sum_ct_c=1$.  The
map sending a design with weight vector $t$ to its chart $P$ is a
bijection from the set of such designs modulo the action
$g_c\mapsto Ug_c$ of the unitary group of $\FF^k$ onto the set of
orthogonal projections of rank $k$ on $\FF^m$.
\end{lemma}

\begin{proof}
Let $P$ be an orthogonal projection of rank $k$ and let $V$ be the
$m\times k$ matrix whose columns are an orthonormal basis of the range of
$P$; then $V\adj V=I_k$ and $VV\adj=P$.  Writing $v_c\adj$ for the $c$-th
row of $V$ and setting $g_c:=t_c^{-1/2}v_c$ produces a family with
$\sum_ct_cg_c^{\vphantom{\mathsf{H}}}g_c\adj=\sum_cv_c^{\vphantom{\mathsf{H}}}v_c\adj
=V\adj V=I_k$, a design with weight vector $t$ and chart $P$.  If two
designs with weight vector $t$ have the same chart, their matrices $V$
and $V'$ have the same column space and orthonormal columns, so their
columns are two orthonormal bases of one $k$-dimensional space and
$V'=VU$ for a unitary $U$ of $\FF^k$, whence $g'_c=U\adj g_c$; conversely
replacing $g_c$ by $Ug_c$ replaces $V$ by $VU\adj$ and leaves
$P=VV\adj$ unchanged.
\end{proof}

\begin{remark}\label{rem:frames}
In frame language \eqref{eq:isometry} says that $(v_c)_{c\le m}$ is a
Parseval frame for $\FF^k$ and $P$ is the associated projection of the
ambient $\FF^m$ onto it \cite{CasazzaKutyniok2013}.  The classical
problem is the uniform slice: for $t_c\equiv1/n$ and $g_c:=\sqrt n\,a_c$,
where $a_c\adj$ is the $c$-th row of a matrix $A$ with $A\adj A=I_k$, the
factor $\sqrt n$ in the vectors cancels the factor $1/n$ in the weights
and $P=AA\adj$; domination of $C$ reads
$\sum_{c\in C}a_c^{\vphantom{\mathsf{H}}}a_c\adj\psd n^{-1}I_k$, which is the
row-selection problem of \cite{GTZ1997}.  In
that slice $D=n^{-1}I$ and the chart criterion below is the condition
$\lmin(P[C])\ge1/n$ used in \cite{Nesterenko2023,Nesterenko2025}.  The
general weight vector is new: $D$ is then a genuine diagonal, and the
comparison is Loewner rather than scalar.
\end{remark}

\subsection{The criterion}\label{sec:criterion}

Two steps separate $\SC$ from the block $W[C]$: a congruence by a
positive diagonal, which holds at every selection size, and an exchange
of the two products in \eqref{eq:twoproducts}, which needs
$\lvert C\rvert=k$.

\begin{lemma}[Lean]\label{lem:congruence}
For every index set $C$ of size $r$,
\begin{equation}\label{eq:congruence}
  W[C]\;=\;\diag(\sqrt{t_c})_{c\in C}\;
  \bigl(\Gram_C-I_r\bigr)\;\diag(\sqrt{t_c})_{c\in C},
\end{equation}
and therefore $W[C]\psd0$ if and only if $\Gram_C\psd I_r$.
\end{lemma}

\begin{proof}
Entrywise, for $c,d\in C$,
\[
  W[C]_{cd}
  \;=\;P_{cd}-t_c\delta_{cd}
  \;=\;\sqrt{t_ct_d}\,\langle g_c,g_d\rangle-\sqrt{t_ct_d}\,\delta_{cd}
  \;=\;\sqrt{t_ct_d}\,\bigl(\Gram_C-I_r\bigr)_{cd},
\]
where the middle step uses $\sqrt{t_ct_d}\,\delta_{cd}=t_c\delta_{cd}$.
This is \eqref{eq:congruence}.  The matrix $B:=\diag(\sqrt{t_c})_{c\in C}$
is invertible because every $t_c>0$, and for invertible $B$ one has
$X\psd0$ if and only if $B\adj XB\psd0$: the substitutions $x\mapsto Bx$
and $x\mapsto B^{-1}x$ carry each quadratic form to the other.
\end{proof}

\begin{lemma}[Lean]\label{lem:flip}
Let $X$ be a matrix over $\FF$, of any shape.  Then
\[
  I-X\adj X\psd0
  \quad\Longleftrightarrow\quad
  I-XX\adj\psd0 .
\]
If $M$ is square, then $M\adj M\psd I$ if and only if $MM\adj\psd I$.
\end{lemma}

\begin{proof}
Both conditions say that a map does not increase length: $I-X\adj X\psd0$
says $\lVert Xu\rVert\le\lVert u\rVert$ for every $u$, and $I-XX\adj\psd0$
says $\lVert X\adj w\rVert\le\lVert w\rVert$ for every $w$.  The first
statement therefore asserts that $X$ is a contraction if and only if
$X\adj$ is, and being symmetric in the two it needs only one implication.

Let $X$ be a contraction and fix $w$.  Then
$\lVert X\adj w\rVert^{2}=\langle w,XX\adj w\rangle$; applying
Cauchy--Schwarz to that inner product and then the hypothesis at
$u=X\adj w$,
\[
  \lVert X\adj w\rVert^4
  \;=\;\langle w,X(X\adj w)\rangle^2
  \;\le\;\lVert w\rVert^2\,\lVert X(X\adj w)\rVert^2
  \;\le\;\lVert w\rVert^2\,\lVert X\adj w\rVert^2 .
\]
If $X\adj w=0$ the bound $\lVert X\adj w\rVert\le\lVert w\rVert$ is
immediate; otherwise divide by $\lVert X\adj w\rVert^{2}$.

The second statement carries the opposite inequality: $M\adj M\psd I$ says
$\lVert Mu\rVert\ge\lVert u\rVert$, so $M$ expands rather than contracts.
Inverting reduces it to the case just proved.  An expansion is injective,
so the square matrix $M$ is invertible, and substituting $u=M^{-1}w$ gives
$\lVert w\rVert\ge\lVert M^{-1}w\rVert$: the inverse is a contraction.  By
the first statement its adjoint $(M^{-1})\adj=(M\adj)^{-1}$ is a
contraction as well, and
substituting $w=M\adj u$ in
$\lVert (M\adj)^{-1}w\rVert\le\lVert w\rVert$ returns
$\lVert u\rVert\le\lVert M\adj u\rVert$, which is $MM\adj\psd I$.  The
converse exchanges $M$ and $M\adj$.
\end{proof}

\begin{lemma}[Lean]\label{lem:chart}
For every $k$-subset $C$,
\[
  \SC\psd I_k \quad\Longleftrightarrow\quad W[C]\psd 0 .
\]
\end{lemma}

\begin{proof}
For $\lvert C\rvert=k$ the matrix $M_C$ is square, so
Lemma~\ref{lem:flip} applies to \eqref{eq:twoproducts} and gives
$\SC=M_C\adj M_C^{\vphantom{\mathsf{H}}}\psd I_k$ if and only if
$\Gram_C=M_C^{\vphantom{\mathsf{H}}}M_C\adj\psd I_k$.  By
Lemma~\ref{lem:congruence} the latter is $W[C]\psd0$.
\end{proof}

The problem therefore reads: \emph{some $k\times k$ principal block of a
rank-$k$ orthogonal projection dominates the corresponding block of a
positive diagonal matrix of trace one}.

The weights enter only through $D$, so the projection and the weights may
be varied independently, and $W[C]\psd0$ is a closed condition in both
arguments.  The compactification of \S\ref{sec:compactness} rests on both
facts.  Figure~\ref{fig:chart} traces the construction.

\begin{figure}[ht]
  \centering
  \begin{tikzpicture}[
  x=3.72mm, y=3.72mm,
  every node/.style={inner sep=1.5pt},
  cellline/.style={draw=black!45, line width=0.25pt},
  border/.style  ={draw=black!85, line width=0.5pt},
  ar/.style      ={-{Latex[length=2mm,width=1.5mm]}, line width=0.5pt}]

% a grid of #3 columns and #4 rows with lower left corner (#1,#2)
\newcommand{\gridframe}[4]{%
  \foreach \gi in {0,...,#3}{\draw[cellline] (#1+\gi,#2) -- (#1+\gi,#2+#4);}
  \foreach \gj in {0,...,#4}{\draw[cellline] (#1,#2+\gj) -- (#1+#3,#2+\gj);}
  \draw[border] (#1,#2) rectangle (#1+#3,#2+#4);}
% cell in row #2 from the top and column #3 of the six by six matrix whose
% lower left corner is (#1,-3), filled with #4
\newcommand{\wfill}[4]{\path[fill=#4] (#1+#3-1,3-#2) rectangle ++(1,1);}

% -------------------------- the design -------------------------------
\node[align=left, font=\small] at (3.2,0)
  {$g_1,\dots,g_m\in\FF^{k}$\\[3pt]
   $t_c>0,\quad\textstyle\sum_c t_c=1$\\[3pt]
   $\textstyle\sum_c t_c\,g_c^{\vphantom{*}}g_c^{*}=I_k$};

\draw[ar] (7.0,0) -- node[above, font=\scriptsize] {$v_c=\sqrt{t_c}\,g_c$} (10.3,0);

% ----------------------------- V -------------------------------------
\gridframe{10.5}{-3}{3}{6}
\node[font=\scriptsize, anchor=south] at (12.0,3.4) {$V^{*}V=I_k$};
\node[font=\scriptsize, align=center, anchor=north] at (12.0,-3.4)
  {$V$, rows $v_c^{*}$\\ $m\times k$};

\draw[ar] (13.9,0) -- node[above, font=\scriptsize] {$P=VV^{*}$} (17.0,0);

% ----------------------------- P -------------------------------------
\gridframe{17.2}{-3}{6}{6}
\node[font=\scriptsize, anchor=south] at (20.2,3.4) {$P^{*}=P=P^{2}$, \ $\rk P=k$};
\node[font=\scriptsize, anchor=north] at (20.2,-3.4)
  {$P_{cd}=\sqrt{t_ct_d}\,\langle g_c,g_d\rangle$};

\draw[ar] (23.5,0) -- (28.7,0);
\node[font=\scriptsize] at (26.1,1.05) {$-\,D$};
\node[font=\scriptsize] at (26.1,-1.05) {$D=\diag(t)$};

% ----------------------------- W -------------------------------------
% the selected rows and columns are C = {1,3,4}
\foreach \r in {1,3,4}{\foreach \c in {1,...,6}{\wfill{28.9}{\r}{\c}{black!12}}}
\foreach \c in {1,3,4}{\foreach \r in {1,...,6}{\wfill{28.9}{\r}{\c}{black!12}}}
\foreach \r in {1,3,4}{\foreach \c in {1,3,4}{\wfill{28.9}{\r}{\c}{black!38}}}
\gridframe{28.9}{-3}{6}{6}
\foreach \c in {1,...,6}{\node[font=\tiny] at (28.4+\c,3.45) {\c};}
\foreach \r in {1,...,6}{\node[font=\tiny] at (28.55,3.5-\r) {\r};}
\node[font=\scriptsize, align=center, anchor=south] at (31.9,4.05)
  {$W=P-D$\\ $W_{cc}=t_c(\ell_c-1)$};

% -------------------- the selection window extracted ------------------
\draw[ar] (31.9,-3.3) -- (31.9,-4.1);
\foreach \r in {1,2,3}{\foreach \c in {1,2,3}{%
  \path[fill=black!38] (30.4+\c-1,-7.2+3-\r) rectangle ++(1,1);}}
\gridframe{30.4}{-7.2}{3}{3}
\node[font=\scriptsize, anchor=east] at (30.1,-5.7) {$W[C]$};

\node[font=\small] at (18.0,-5.7) {$W[C]\psd0\iff\SC\psd I_k$};
\node[font=\scriptsize] at (5.8,-5.7) {$m=6$, \ $k=3$, \ $C=\{1,3,4\}$};

\end{tikzpicture}

  \caption{The chart, at $m=6$, $k=3$ and $C=\{1,3,4\}$.  The design data
  pass through the isometry $V$ of \eqref{eq:isometry} to the projection
  $P=VV\adj$ of \eqref{eq:chart} and the gap $W=P-D$ of
  \eqref{eq:chartgap}.  Selecting $C$ is taking a principal block, and
  Lemma~\ref{lem:chart} identifies its semidefiniteness with domination.}
  \label{fig:chart}
\end{figure}

\begin{remark}\label{rem:size}
The equivalence of Lemma~\ref{lem:chart} depends on the hypothesis
$\lvert C\rvert=k$.  Domination is the $k\times k$ statement
$M_C\adj M_C\psd I_k$, while the block condition is the
$r\times r$ statement $M_C^{\vphantom{\mathsf{H}}}M_C\adj\psd I_r$ with
$r=\lvert C\rvert$.  For $r<k$ the left-hand matrix has rank at most
$r<k$ and can never dominate $I_k$, while the block condition is
unconstrained; for $r>k$ the block matrix has rank at most $k<r$ and
never dominates $I_r$, while $\SC$ may.  Both failures are rank effects,
and $r=k$ is the only size at which the equivalence has content in both
directions.
\end{remark}

Taking determinants in \eqref{eq:congruence} converts the criterion into
an exact scalar test, one principal minor at a time.

\begin{corollary}[Lean]\label{cor:minor}
For every index set $C$,
\[
  \det W[C]\;=\;\Bigl(\prod_{c\in C}t_c\Bigr)\,\det(\Gram_C-I).
\]
\end{corollary}

The factor is the square of $\det\diag(\sqrt{t_c})_{c\in C}$ and is
strictly positive, so $W[C]$ and $\Gram_C-I$ have the same principal
minors up to positive factors and are semidefinite together, minor by
minor.  Over the rationals the minors decide the question by exact
arithmetic, with no eigenvalue computation and no genericity assumption.

The design below has four dominating pairs of six, one of them with slack
zero.

\begin{example}\label{ex:chartworked}
Over $\R$ at $(m,k)=(4,2)$ take
\[
  g_0=(\sqrt2,0),\qquad
  g_1=\tfrac1{\sqrt3}(2,1),\qquad
  g_2=\tfrac1{\sqrt3}(1,-2),\qquad
  g_3=(0,2)
\]
with weights $t=\tfrac19(2,3,3,1)$, which are positive and sum to one.
The four terms of \eqref{eq:parseval} are
\[
  \tfrac19\begin{pmatrix}4&0\\0&0\end{pmatrix},\qquad
  \tfrac19\begin{pmatrix}4&2\\2&1\end{pmatrix},\qquad
  \tfrac19\begin{pmatrix}1&-2\\-2&4\end{pmatrix},\qquad
  \tfrac19\begin{pmatrix}0&0\\0&4\end{pmatrix},
\]
of sum $\tfrac19\begin{psmallmatrix}9&0\\0&9\end{psmallmatrix}=I_2$: a
design at $(4,2)$.  The leverages are
$\ell=(2,\tfrac53,\tfrac53,4)$ and the leverage scores
$s_c=t_c\ell_c=\tfrac19(4,5,5,4)$, which sum to $2=k$
(Lemma~\ref{lem:trace}).

Whitening clears the radicals.  With $v_c=\sqrt{t_c}\,g_c$,
\[
  V=\tfrac13\begin{pmatrix}2&0\\2&1\\1&-2\\0&2\end{pmatrix},
\]
and the two columns of $3V$ have squared lengths $4+4+1+0=9$ and
$0+1+4+4=9$ and inner product $0+2-2+0=0$.  Dividing by $9$ gives
$V\adj V=I_2$, which is \eqref{eq:isometry}.  The atoms carry $\sqrt2$ and
$\sqrt3$ and the chart carries neither:
\[
  P=VV\adj=\tfrac19\begin{pmatrix}
    4&4&2&0\\ 4&5&0&2\\ 2&0&5&-4\\ 0&2&-4&4\end{pmatrix}.
\]
Idempotency is ten integer identities.  Write $r_c$ for the $c$-th row of
$9P$.  Then $\bigl((9P)^2\bigr)_{cd}=\langle r_c,r_d\rangle$, and $P^2=P$
asks that this equal $9\,(9P)_{cd}$.  The ten inner products are
\[
  \begin{array}{llll}
  \langle r_0,r_0\rangle=36, & \langle r_0,r_1\rangle=36,
    & \langle r_0,r_2\rangle=18, & \langle r_0,r_3\rangle=0,\\[2pt]
  \langle r_1,r_1\rangle=45, & \langle r_1,r_2\rangle=0,
    & \langle r_1,r_3\rangle=18, & \langle r_2,r_2\rangle=45,\\[2pt]
  \langle r_2,r_3\rangle=-36, & \langle r_3,r_3\rangle=36,
  \end{array}
\]
which are $9$ times $4,4,2,0,5,0,2,5,-4,4$, the corresponding entries of
$9P$.  The trace is $\tfrac19(4+5+5+4)=2=k$ and the diagonal is $s$, as
Proposition~\ref{prop:chartid} says.  Subtracting $D=\diag(t)$,
\[
  W=\tfrac19\begin{pmatrix}
    2&4&2&0\\ 4&2&0&2\\ 2&0&2&-4\\ 0&2&-4&3\end{pmatrix},
  \qquad
  W_{cc}=t_c(\ell_c-1)=\tfrac19(2,2,2,3),
  \qquad \textstyle\sum_cW_{cc}=1=k-1 .
\]

Now take the six blocks.  Every $W_{cc}$ is positive, so a $2\times2$
principal block of $W$ is positive semidefinite precisely when its
determinant is nonnegative and positive definite precisely when that
determinant is positive:
\[
  \begin{array}{r|rrrrrr}
  C & \{0,1\} & \{0,2\} & \{0,3\} & \{1,2\} & \{1,3\} & \{2,3\}\\[2pt]
  81\det W[C] & -12 & 0 & 6 & 4 & 2 & -10
  \end{array}
\]
By Lemma~\ref{lem:chart} four pairs dominate, three of them strictly,
while $\{0,1\}$ and $\{2,3\}$ do not.

Corollary~\ref{cor:minor} checks each entry a second way.  For a pair,
$\det(\Gram_C-I)=(\ell_c-1)(\ell_d-1)-\lvert\langle
g_c,g_d\rangle\rvert^2$, in which the inner product enters squared, so the
radicals cancel there too.  Here $\ell_c-1=(1,\tfrac23,\tfrac23,3)$, and
the six squared inner products, in the order of the table below, are
$\tfrac83$, $\tfrac23$, $0$, $0$, $\tfrac43$, $\tfrac{16}3$, whence
\[
  \begin{array}{r|rrrrrr}
  C & \{0,1\} & \{0,2\} & \{0,3\} & \{1,2\} & \{1,3\} & \{2,3\}\\[2pt]
  \prod_{c\in C}t_c
    & \tfrac2{27} & \tfrac2{27} & \tfrac2{81}
    & \tfrac19 & \tfrac1{27} & \tfrac1{27}\\[4pt]
  \det(\Gram_C-I) & -2 & 0 & 3 & \tfrac49 & \tfrac23 & -\tfrac{10}3
  \end{array}
\]
The six products of the last two rows, multiplied by $81$, return the
minors of the first table.  At $\{0,3\}$ the product is
$\tfrac2{81}\cdot3=\tfrac6{81}$ and $81\cdot\tfrac6{81}=6$.

The pair with determinant zero is $\{0,2\}$, where
\[
  \SC=\begin{pmatrix}2&0\\0&0\end{pmatrix}
     +\tfrac13\begin{pmatrix}1&-2\\-2&4\end{pmatrix}
     =\begin{pmatrix}\tfrac73&-\tfrac23\\[3pt]
                     -\tfrac23&\tfrac43\end{pmatrix} .
\]
The matrix $\SC-I_2$ has determinant $0$ and trace $\tfrac53$, so
$\operatorname{spec}\SC=\{1,\tfrac83\}$ and $\lmin(\SC)=1$.  Its kernel is
spanned by $(1,2)$, and at $x=\tfrac1{\sqrt5}(1,2)$ the pointwise form
\eqref{eq:pointwise} reads $\tfrac25+\tfrac35=1$: the pair covers that one
direction with nothing to spare and every other direction strictly.

The two failures need no determinant.  A dominating pair has
$(\SC)_{jj}\ge1$ at every $j$, and
\[
  S_{\{0,1\}}=\begin{pmatrix}\tfrac{10}3&\tfrac23\\[3pt]
                             \tfrac23&\tfrac13\end{pmatrix},
  \qquad
  S_{\{2,3\}}=\begin{pmatrix}\tfrac13&-\tfrac23\\[3pt]
                             -\tfrac23&\tfrac{16}3\end{pmatrix}
\]
fall to $\tfrac13$ at $e_2$ and at $e_1$.  On the chart side that shortcut
is unavailable, since every $W_{cc}$ is positive.  Over the six pairs
$\max_C\lmin(\SC)=2$, attained at $\{0,3\}$, where $\SC=\diag(2,4)$.
\end{example}

The six pairs are the edges of $K_4$ on the four indices.
Figure~\ref{fig:chartworked} labels each edge with its minor, and the two
failing pairs come out opposite edges.

\begin{figure}[ht]
  \centering
  % The chart of one real design at (4,2), carried through and read off at all
% six pairs.  Three panels: the pipeline g -> V -> P -> W with the matrices in
% it, the four atoms against the unit circle of F^2 = R^2, and the six pairs as
% the edges of K_4 labelled by 81 det W[C].
%
% PROJECTION.  The atom panel is a plane figure and carries no projection of a
% higher-dimensional space: the map is the identity of R^2 scaled by 1.05cm,
%   X = 1.05 x_1,   Y = 1.05 x_2,
% so no viewing normal and no visibility rule are needed and no line is hidden.
% Every coordinate in that panel is that pair.  The K_4 panel is a graph
% drawing on the square derived below, and the reading block is a stack of
% rows.  Every coordinate in all three carries its exact source in a trailing
% comment.
%
% THE DESIGN, exactly.  Over R, m = 4, k = 2.
%   g_0 = (sqrt2, 0)        t_0 = 2/9    ell_0 = 2      s_0 = 4/9
%   g_1 = (2, 1)/sqrt3      t_1 = 3/9    ell_1 = 5/3    s_1 = 5/9
%   g_2 = (1,-2)/sqrt3      t_2 = 3/9    ell_2 = 5/3    s_2 = 5/9
%   g_3 = (0, 2)            t_3 = 1/9    ell_3 = 4      s_3 = 4/9
% The weights are positive and sum to one, and sum_c t_c g_c g_c^T = I_2 over
% the rationals: the four rank-one terms are (2/9)[[2,0],[0,0]],
% (1/9)[[4,2],[2,1]], (1/9)[[1,-2],[-2,4]] and (1/9)[[0,0],[0,4]], of sum
% (1/9)[[9,0],[0,9]].  The leverage scores sum to 4/9+5/9+5/9+4/9 = 2 = k.
% Whitening gives v_c = sqrt(t_c) g_c = (1/3)(2,0), (1/3)(2,1), (1/3)(1,-2),
% (1/3)(0,2), whence the rational V, P = V V^T and W = P - diag(t) printed in
% the pipeline.  81 det W[C] over the six pairs is -12, 0, +6, +4, +2, -10 in
% the order {0,1}, {0,2}, {0,3}, {1,2}, {1,3}, {2,3}.
%
% ELLIPSE FRAME.  For a pair C the set {x : x^T S_C x = 1} is an ellipse with
% semi-axes 1/sqrt(lambda) along the eigenvectors of S_C, and S_C >= I holds
% exactly when that ellipse lies in the closed unit disc.  At C = {0,2},
% S_C = [[7/3,-2/3],[-2/3,4/3]] has spectrum {1, 8/3}, so the semi-axes are 1
% and sqrt(3/8) = 0.6123724, the long one along the eigenvector (1,2)/sqrt5 at
% atan2(2,1) = 63.4349 degrees, and the ellipse touches the unit circle at the
% two points +-(1,2)/sqrt5 = +-(0.4472136, 0.8944272).  Only this ellipse is
% drawn.  The two failing pairs have long semi-axis 2.0163 at C = {2,3} and
% 2.2830 at C = {0,1}, and bounding half-extents sqrt((S_C^-1)_jj) of (2, 1/2)
% and (1/sqrt2, sqrt5) = (0.7071068, 2.2360680) against the 1.05 of the unit
% circle, so each widens the plate in one direction; and the K_4 panel already
% records the sign that decides them.
%
% K_4 FRAME.  The four indices sit at the corners of a square of half-side
% 1.05cm in the cyclic order 0, 1, 2, 3, so that the complement map C -> C^c on
% pairs is the opposite-edge map: {0,1} against {2,3}, {0,3} against {1,2}, and
% the two diagonals {0,2} against {1,3}.  That is the pairing thm:duality uses
% at this self-dual cell; figures/dualworked.tex runs the same pairing down its
% table, C in the first column against E = C^c in the third.  Edge labels sit
% at parameter 0.28 from the first endpoint on the two diagonals and at the
% midpoint on the four sides; the open marker on {0,2} sits at parameter 0.75
% along it, clear of the crossing.  The tangency label names the line through
% the two marks and so carries the sign +-, as the caption does.
%
% DISCLOSED: this design is new to the paper and no Lean declaration mentions
% it.  Every number above and below was computed in exact arithmetic outside
% the assistant -- rational for the chart, the gap and the six minors,
% algebraic for the atoms, the ellipse and the tangency, whose decimals in the
% body are truncations of those exact values.  What is mechanized is the
% machinery it instantiates,
% not the instance: the chart identities of prop:chartid, the criterion of
% lem:chart, and the minor identity of cor:minor.
\begin{tikzpicture}[
  x=1cm, y=1cm,
  ar/.style     ={-{Latex[length=2mm,width=1.5mm]}, line width=0.5pt},
  vec/.style    ={-{Latex[length=2mm,width=1.5mm]}, line width=0.75pt},
  rim/.style    ={line width=0.6pt, black!70},
  lev/.style    ={line width=0.5pt, black!70},
  frame/.style  ={densely dotted, line width=0.3pt, black!35},
  tight/.style  ={line width=0.5pt, black!55, dashed,
                  dash pattern=on 2.2pt off 1.6pt},
  win/.style    ={line width=1.1pt},
  lose/.style   ={line width=0.55pt, black!55, dashed,
                  dash pattern=on 1.4pt off 1.2pt},
  tag/.style    ={font=\footnotesize, fill=white, inner sep=0.8pt},
  blk/.style    ={inner sep=1.5pt, font=\scriptsize},
  every node/.style={inner sep=1.5pt, font=\small}]

% =========================== the pipeline ================================
\node[blk, anchor=west, align=left] (A) at (0,4.75)
  {$m=4$, \ $k=2$, \ over $\R$\\[3pt]
   $g_0=(\sqrt2,0)$\\[2pt]
   $g_1=\tfrac1{\sqrt3}(2,1)$\\[2pt]
   $g_2=\tfrac1{\sqrt3}(1,-2)$\\[2pt]
   $g_3=(0,2)$\\[3pt]
   $t=\tfrac19(2,3,3,1)$};
\node[blk, right=15mm of A] (Vm)
  {$V=\tfrac13\begin{psmallmatrix}2&0\\2&1\\1&-2\\0&2\end{psmallmatrix}$};
\node[blk, right=15mm of Vm] (Pm)
  {$P=\tfrac19\begin{psmallmatrix}
      4&4&2&0\\4&5&0&2\\2&0&5&-4\\0&2&-4&4\end{psmallmatrix}$};
\node[blk, right=15mm of Pm] (Wm)
  {$W=\tfrac19\begin{psmallmatrix}
      2&4&2&0\\4&2&0&2\\2&0&2&-4\\0&2&-4&3\end{psmallmatrix}$};

\draw[ar] (A)  -- node[above, font=\scriptsize] {$\sqrt{t_c}\,g_c$} (Vm);
\draw[ar] (Vm) -- node[above, font=\scriptsize] {$VV\adj$} (Pm);
\draw[ar] (Pm) -- node[above, font=\scriptsize] {$-\,D$} (Wm);

\node[blk, below=2mm of Vm] {$V\adj V=I_2$};
\node[blk, below=2mm of Pm, align=center]
  {$P^2=P$, \ $\tr P=2$\\[2pt]
   $s=\tfrac19(4,5,5,4)$, \ $\ell=(2,\tfrac53,\tfrac53,4)$};
\node[blk, below=2mm of Wm, align=center]
  {$W_{cc}=t_c(\ell_c-1)$\\[2pt] $\textstyle\sum_cW_{cc}=k-1=1$};

% ======================= the atoms, against the unit circle ================
\begin{scope}[shift={(2.05,0)}]
  \draw[frame] (-1.75,0)--(1.78,0);
  \draw[frame] (0,-1.62)--(0,2.42);
  \node[anchor=west,  font=\footnotesize] at ( 1.82, 0.00) {$e_1$};
  \node[anchor=south, font=\footnotesize] at ( 0.00, 2.46) {$e_2$};

  \draw[rim] (0,0) circle (1.05);
  % the ellipse x^T S_C x = 1 at C = {0,2}: semi-axes 1 and sqrt(3/8), the
  % long one along (1,2)/sqrt5
  \draw[lev, rotate=63.4349] (0,0) ellipse (1.05 and 0.64299);
  \draw[tight] (-0.58227,-1.16454)--(0.58227,1.16454);       % +-1.24 (1,2)/sqrt5
  \draw[lev, fill=white] ( 0.46957, 0.93915) circle (1.4pt); % + (1,2)/sqrt5
  \draw[lev, fill=white] (-0.46957,-0.93915) circle (1.4pt); % - (1,2)/sqrt5

  \draw[vec] (0,0)--( 1.48492, 0.00000);   % g_0 = (sqrt2, 0)
  \draw[vec] (0,0)--( 1.21244, 0.60622);   % g_1 = (2, 1)/sqrt3
  \draw[vec] (0,0)--( 0.60622,-1.21244);   % g_2 = (1,-2)/sqrt3
  \draw[vec] (0,0)--( 0.00000, 2.10000);   % g_3 = (0, 2)
  \fill (0,0) circle (1.1pt);
  \node[tag, anchor=north]      at ( 1.46000,-0.09) {$g_0$};
  \node[tag, anchor=south west] at ( 1.23000, 0.63) {$g_1$};
  \node[tag, anchor=west]       at ( 0.65000,-1.24) {$g_2$};
  \node[tag, anchor=east]       at (-0.07000, 1.92) {$g_3$};
  % names the tangency LINE, whose two ends are +-(1,2)/sqrt5; set clear of the
  % dashed segment, which stops at x = -0.58227
  \node[tag, anchor=north east] at (-0.62000,-1.02) {$\pm\tfrac1{\sqrt5}(1,2)$};

  \node[anchor=north west, font=\footnotesize] at (-1.75,-1.72)
    {$x\adj\SC x=1$ at $C=\{0,2\}$, \ $\operatorname{spec}\SC=\{1,\tfrac83\}$};
\end{scope}

% ========================= the six pairs as K_4 ============================
\begin{scope}[shift={(6.65,0.40)}]
  \coordinate (n0) at (-1.05, 1.05);
  \coordinate (n1) at ( 1.05, 1.05);
  \coordinate (n2) at ( 1.05,-1.05);
  \coordinate (n3) at (-1.05,-1.05);

  \draw[lose] (n0)--(n1);          % {0,1}
  \draw[lose] (n2)--(n3);          % {2,3}
  \draw[win]  (n0)--(n3);          % {0,3}
  \draw[win]  (n1)--(n2);          % {1,2}
  \draw[win]  (n0)--(n2);          % {0,2}
  \draw[win]  (n1)--(n3);          % {1,3}
  \draw[lev, fill=white] (0.52500,-0.52500) circle (1.4pt);  % 0.75 along {0,2}

  \foreach \nn in {n0,n1,n2,n3}{\fill (\nn) circle (1.5pt);}
  \node[anchor=south east, font=\footnotesize] at (n0) {$0$};
  \node[anchor=south west, font=\footnotesize] at (n1) {$1$};
  \node[anchor=north west, font=\footnotesize] at (n2) {$2$};
  \node[anchor=north east, font=\footnotesize] at (n3) {$3$};

  \node[tag] at ( 0.00000, 1.05000) {$-12$};   % midpoint {0,1}
  \node[tag] at ( 0.00000,-1.05000) {$-10$};   % midpoint {2,3}
  \node[tag] at (-1.05000, 0.00000) {$+6$};    % midpoint {0,3}
  \node[tag] at ( 1.05000, 0.00000) {$+4$};    % midpoint {1,2}
  \node[tag] at (-0.46200, 0.46200) {$0$};     % 0.28 along {0,2}
  \node[tag] at ( 0.46200, 0.46200) {$+2$};    % 0.28 along {1,3}

  \node[anchor=north, font=\footnotesize] at (0,-1.44) {$81\det W[C]$};
\end{scope}

% ============================ the readings ================================
\begin{scope}[every node/.style={anchor=west, inner sep=1.5pt,
                                 font=\footnotesize}]
  \draw[win]  (9.30, 1.30)--(9.85, 1.30);
  \node at (9.98, 1.30) {$\det W[C]>0$: dominates strictly};
  \draw[win]  (9.30, 0.86)--(9.85, 0.86);
  \draw[lev, fill=white] (9.575, 0.86) circle (1.4pt);
  \node at (9.98, 0.86) {$=0$: dominates, no slack};
  \draw[lose] (9.30, 0.42)--(9.85, 0.42);
  \node at (9.98, 0.42) {$<0$: fails};

  \node at (9.30,-0.18) {$\lmin(S_{\{0,2\}})=1$, attained at
    $\tfrac1{\sqrt5}(1,2)$};
  \node at (9.30,-0.66) {$(S_{\{0,1\}})_{22}=(S_{\{2,3\}})_{11}=\tfrac13<1$};
  \node at (9.30,-1.14) {$W_{cc}>0$ at all four indices};
\end{scope}

\end{tikzpicture}

  \caption{The design of Example~\ref{ex:chartworked} at $(4,2)$.  Top: the
  atoms pass through $V$ to the chart $P=VV\adj$ of \eqref{eq:chart} and the
  gap $W=P-D$ of \eqref{eq:chartgap}; the atoms carry $\sqrt2$ and $\sqrt3$
  and $9P$ is integral.  Left: at $C=\{0,2\}$ the ellipse $x\adj\SC x=1$ lies
  in the closed unit disc, which for a positive definite $\SC$ is what
  $\SC\psd I_2$ says, and touches the circle at $\pm\tfrac1{\sqrt5}(1,2)$,
  where the slack vanishes.  Centre: the six pairs as the edges of $K_4$,
  labelled by $81\det W[C]$; the complement
  map $C\mapsto C^{\complement}$ of \S\ref{sec:duality} carries each edge to
  the edge opposite it.  Right: the edge styles, and the diagonal entry
  $\tfrac13$ that settles each failure on the design side while every
  $W_{cc}>0$ leaves the chart side needing the minor.  No Lean declaration
  mentions this design.}
  \label{fig:chartworked}
\end{figure}

\subsection{Inertia, and what it does not decide}\label{sec:inertia}

The gap $W$ has a fixed signature, which by itself decides nothing.

\begin{lemma}[Lean]\label{lem:rayleigh}
Let $m\ge2$.  If $Px=x$ then
$\langle x,Wx\rangle=\sum_c(1-t_c)\lvert x_c\rvert^2$, which is positive
for $x\ne0$.  If $Px=0$ then
$\langle x,Wx\rangle=-\sum_ct_c\lvert x_c\rvert^2$, which is negative for
$x\ne0$.
\end{lemma}

\begin{proof}
$\langle x,Dx\rangle=\sum_ct_c\lvert x_c\rvert^2$ in both cases.  If
$Px=x$ then $\langle x,Px\rangle=\lVert x\rVert^2=\sum_c\lvert
x_c\rvert^2$, and subtracting gives the first identity; each coefficient
$1-t_c$ is strictly positive because $m\ge2$ forces $t_c<1$.  If $Px=0$
then $\langle x,Px\rangle=0$ and the second identity follows; each $t_c$
is strictly positive.
\end{proof}

\begin{lemma}\label{lem:inertia}
For $m\ge2$ the matrix $W$ has inertia $(k,\,m-k,\,0)$.
\end{lemma}

\begin{proof}
The subspaces $\operatorname{ran}P$ and $\ker P$ are complementary of
dimensions $k$ and $m-k$, and by Lemma~\ref{lem:rayleigh} the form
$\langle x,Wx\rangle$ is positive definite on the first and negative
definite on the second.  Sylvester's law of inertia
\cite[\S4.5]{HornJohnson2013} gives the count.
\end{proof}

The hypothesis $m\ge2$ is necessary: at $m=k=1$ the only design has
$t_1=1$ and $\ell_1=1$, so $P=D=1$ and $W=0$, of inertia $(0,0,1)$.

\begin{remark}[Lean]\label{rem:nogo}
A positive semidefinite matrix has a nonnegative diagonal.  Hence a
Hermitian matrix whose diagonal entries are all negative has no positive
semidefinite principal submatrix of any nonempty size.  At $(m,k)=(2,1)$
take
\[
  B_2=\begin{pmatrix}-1&3\\3&-1\end{pmatrix},
  \qquad \det B_2=1-9=-8 .
\]
Its form is $\langle x,B_2x\rangle=-x_1^2-x_2^2+6x_1x_2$, which takes the
value $4$ at $(1,1)$ and $-8$ at $(1,-1)$: the inertia is $(1,1,0)$,
which by Lemma~\ref{lem:inertia} is the inertia of a chart gap at
$(2,1)$.  At $(m,k)=(3,2)$ take
\[
  B_3=\begin{pmatrix}-1&2&2\\2&-1&-2\\2&-2&-1\end{pmatrix},
  \qquad \det B_3=-5 .
\]
On the plane spanned by $(1,1,0)$ and $(1,0,1)$, parameterized as
$x=(a+b,a,b)$,
\[
  \langle x,B_3x\rangle
  \;=\;-(a+b)^2-a^2-b^2+4a(a+b)+4b(a+b)-4ab
  \;=\;2a^2+2ab+2b^2,
\]
which is $(a+b)^2+a^2+b^2$, hence positive definite; so $B_3$ has at
least two positive eigenvalues, and $\det B_3<0$ excludes three, leaving
inertia $(2,1,0)$, the inertia of a chart gap at $(3,2)$.  Every diagonal
entry of $B_2$ and of $B_3$ is $-1$, so neither matrix has a positive
semidefinite principal block of the required size.  No statement of the
form ``a Hermitian matrix of inertia $(k,m-k,0)$ possesses a positive
semidefinite $k\times k$ principal block'' is therefore available, and no
proof of $\Gtz$ can run on the signature of $W$ alone.

A chart gap differs from $B_2$ and $B_3$ in its diagonal: by
Proposition~\ref{prop:chartid} it has $W_{cc}=t_c(\ell_c-1)\ge-t_c$ at
every index, with $-t_c>-1$ once $m\ge2$, and $\sum_cW_{cc}=k-1\ge0$, so
at least one diagonal entry is nonnegative.
\end{remark}

\subsection{Duality}\label{sec:duality}

Since $I-P$ is the orthogonal projection onto $\ker P$, selecting a block
of $P$ and deleting the complementary block of $I-P$ are the same
operation.  With the correct normalization that operation turns a design
of rank $k$ into a design of rank $m-k$, exchanging domination for
domination at the complement.

Fix a design at $(m,k)$ with $k\ge1$ and $m\ge k+1$, and put $r:=m-k\ge1$;
then $m\ge2$ and every $t_c<1$.  Choose an $m\times r$ matrix $Z$ whose
columns are an orthonormal basis of $\ker P$, so that
\[
  Z\adj Z=I_r,\qquad ZZ\adj=I_m-P ,
\]
and write $z_c\adj$ for the $c$-th row of $Z$.  As in
\eqref{eq:isometry} and Proposition~\ref{prop:chartid},
\[
  \sum_c z_c^{\vphantom{\mathsf{H}}}z_c\adj=Z\adj Z=I_r,
  \qquad
  \lVert z_c\rVert^2=(ZZ\adj)_{cc}=1-P_{cc}=1-s_c .
\]
Rescale by the co-weights: put
\[
  w_c\;:=\;\frac{z_c}{\sqrt{1-t_c}}\in\FF^{r},
\]
which is legitimate because $t_c<1$.  The first display becomes
\begin{equation}\label{eq:kernelframe}
  \sum_c(1-t_c)\,w_c^{\vphantom{\mathsf{H}}}w_c\adj\;=\;I_r .
\end{equation}
This is not a Parseval condition, since the coefficients $1-t_c$ sum to
$m-1$ rather than to $1$.  The correct rescaling is by the operator
\begin{equation}\label{eq:dualgram}
  \Sdual\;:=\;\sum_c t_c\,w_c^{\vphantom{\mathsf{H}}}w_c\adj .
\end{equation}

\begin{lemma}\label{lem:dualgram}
$\Sdual\pd0$.
\end{lemma}

\begin{proof}
$\Sdual\psd0$ as a nonnegative combination of rank-one positive semidefinite
matrices.  If $\Sdual u=0$ then $\langle u,\Sdual u\rangle=\sum_ct_c\lvert\langle
w_c,u\rangle\rvert^2=0$, so $\langle w_c,u\rangle=0$ for every $c$
because every $t_c>0$; then \eqref{eq:kernelframe} applied to $u$ gives
$u=\sum_c(1-t_c)\langle w_c,u\rangle\,w_c=0$.
\end{proof}

\begin{lemma}[Lean]\label{lem:whiten}
Every $A\pd0$ admits an invertible $R$ with $R\adj AR=I$.
\end{lemma}

\begin{proof}
Write $A=LL\adj$ with $L$ invertible, by Cholesky factorization, and put
$R:=(L\adj)^{-1}$.  Then $R\adj AR=L^{-1}LL\adj(L\adj)^{-1}=I$.
\end{proof}

\begin{theorem}[Lean]\label{thm:duality}
Let $k\ge1$ and $m\ge k+1$, and let $\dsn$ be a design at $(m,k)$
over $\FF$ with weight vector $t$.  With $w_c$, $\Sdual$ and $R$ as above, the
vectors
\[
  h_c\;:=\;R\adj w_c\in\FF^{m-k}
\]
with the \emph{same} weights $t_c$ form a design $\dsn^{\perp}$ at
$(m,m-k)$, and for every $k$-subset $C$,
\[
  C\ \text{dominates in}\ \dsn
  \quad\Longleftrightarrow\quad
  C^{\complement}\ \text{dominates in}\ \dsn^{\perp} .
\]
\end{theorem}

\begin{proof}
The weights are unchanged, so they are positive and sum to one, and
\[
  \sum_c t_c\,h_c^{\vphantom{\mathsf{H}}}h_c\adj
  \;=\;R\adj\Bigl(\sum_ct_c\,w_c^{\vphantom{\mathsf{H}}}w_c\adj\Bigr)R
  \;=\;R\adj\Sdual R\;=\;I_{m-k}
\]
by \eqref{eq:dualgram} and Lemma~\ref{lem:whiten}; here the conjugation
identity $R\adj(ww\adj)R=(R\adj w)(R\adj w)\adj$ is used termwise.  So
$\dsn^{\perp}$ is a design at $(m,m-k)$.

Fix $C$ with $\lvert C\rvert=k$ and put $E:=C^{\complement}$, of size
$r=m-k$.  Let $Y$ be the $k\times r$ matrix whose rows are $w_c\adj$,
$c\in C$; equivalently
$Y=\diag\bigl((1-t_c)^{-1/2}\bigr)_{c\in C}\,Z_C$, where $Z_C$ is the
submatrix of $Z$ on the rows $C$.  Then, as in \eqref{eq:twoproducts},
\[
  Y^{\vphantom{\mathsf{H}}}Y\adj=\bigl(\langle w_c,w_d\rangle\bigr)_{c,d\in C},
  \qquad
  Y\adj Y^{\vphantom{\mathsf{H}}}=\sum_{c\in C}w_c^{\vphantom{\mathsf{H}}}w_c\adj .
\]
Now run the equivalences.  Since $P=I-ZZ\adj$,
\[
  W[C]\;=\;\diag(1-t_c)_{c\in C}-(ZZ\adj)[C],
\]
and $(ZZ\adj)[C]=Z_C^{\vphantom{\mathsf{H}}}Z_C\adj$.  Conjugating by the
invertible diagonal $\diag\bigl((1-t_c)^{-1/2}\bigr)_{c\in C}$ as in the
proof of Lemma~\ref{lem:congruence},
\[
  W[C]\psd0
  \quad\Longleftrightarrow\quad
  I_k-Y^{\vphantom{\mathsf{H}}}Y\adj\psd0
  \quad\Longleftrightarrow\quad
  I_r-Y\adj Y^{\vphantom{\mathsf{H}}}\psd0 ,
\]
the second step by Lemma~\ref{lem:flip} applied to $X:=Y\adj$.  By
\eqref{eq:kernelframe} the right-hand condition reads
\[
  \sum_{c\in C}w_c^{\vphantom{\mathsf{H}}}w_c\adj
  \;\preceq\;\sum_{c}(1-t_c)\,w_c^{\vphantom{\mathsf{H}}}w_c\adj .
\]
Write $T:=\sum_c w_c^{\vphantom{\mathsf{H}}}w_c\adj$ for the total over all
indices.  Two identities turn that display inside out: the left side is
$T$ less the sum over $E$, since $C$ and $E$ partition the indices, and
the right side is $T-\Sdual$ by \eqref{eq:dualgram}.  Subtracting both
sides from $T$ therefore reverses the inequality into
\[
  \sum_{c\in E}w_c^{\vphantom{\mathsf{H}}}w_c\adj\;\psd\;\Sdual ,
\]
and conjugating by the invertible $R$ turns it into
$\sum_{c\in E}h_c^{\vphantom{\mathsf{H}}}h_c\adj\psd R\adj\Sdual R=I_{m-k}$, which is
domination of $E$ in $\dsn^{\perp}$.  Every step is an
equivalence, and $\lvert E\rvert=m-k$ is the rank of $\dsn^{\perp}$,
so Lemma~\ref{lem:chart} applies at both ends.
\end{proof}

\begin{corollary}[Lean]\label{cor:dualcell}
For $k\ge1$ and $m\ge k+1$,
$\Gtz(m,k)\Leftrightarrow\Gtz(m,m-k)$, and likewise over $\C$.
\end{corollary}

\begin{proof}
If every design at $(m,m-k)$ has a dominating $(m-k)$-subset, apply this
to $\dsn^{\perp}$ and complement the resulting subset.  The
construction is symmetric in the two cells.
\end{proof}

The dual design is not canonical: $Z$ is determined up to a unitary of
$\FF^{r}$ and $R$ up to a unitary factor, and both ambiguities act on the
$h_c$ by a common unitary, which changes no domination.  The
normalization, by contrast, is forced.  Figure~\ref{fig:duality} draws
both normalizations at $(3,1)$.

\begin{remark}\label{rem:naive}
Rescaling the kernel rows by the weights rather than the co-weights, so
that the dual vectors are $u_c:=z_c/\sqrt{t_c}$, also
produces a design at $(m,m-k)$, since $\sum_ct_cu_c^{\vphantom{\mathsf{H}}}u_c\adj
=\sum_cz_c^{\vphantom{\mathsf{H}}}z_c\adj=I_{m-k}$.  It does not satisfy the
conclusion of Theorem~\ref{thm:duality}.  At uniform weights the two
agree: for $t_c\equiv1/m$ one has $\Sdual=\frac{1}{m-1}I$ by
\eqref{eq:kernelframe}, hence $R=\sqrt{m-1}\,I$ and
$h_c=\sqrt{m-1}\,w_c=\sqrt{m}\,z_c=u_c$.  At non-uniform weights they do
not.
\end{remark}

\begin{figure}[ht]
  \centering
  % Duality at (m,k) = (3,1).  The chart is the orthogonal projection onto a line
% of the index space R^3 and I-P is the projection onto the complementary plane,
% so selecting the coordinate 3 upstairs and deleting the complementary pair
% downstairs are one operation.  The design drawn is
%   t = (9,36,4)/49,   g = (2,-1/3,3/2),   v_c = sqrt(t_c) g_c = (6,-2,3)/7,
% whence P = v v^H, s_c = v_c^2 = (36,4,9)/49, and the singleton {3} dominates
% because s_3 - t_3 = 5/49 > 0.  The two rescalings of the kernel row z_3 are
%   correct  w_3 = z_3/sqrt(1-t_3),  ||w_3||^2 = (40/49)/(45/49) = 8/9 <= 1,
%   naive    u_3 = z_3/sqrt(t_3),    chart (Z Z^H, t) = (I-P, t),
% and the naive dual refuses the pair {1,2}: (I-P-D)[{1,2}] = [4,12;12,9]/49 has
% determinant (36-144)/49^2 < 0.
%
% Orthographic projection of R^3 with azimuth 45 and elevation 30 degrees,
% scaled by 1.85cm:
%   X = 1.85(-0.707107 x_1 + 0.707107 x_2)
%   Y = 1.85(-0.353553 x_1 - 0.353553 x_2 + 0.866025 x_3)
% Every coordinate below is the parenthesised value, with its source triple in a
% trailing comment; the 1.85cm is carried by the x and y unit vectors.  The
% kernel plane is drawn as the rectangle spanned by the orthonormal pair
%   d1 =  cos25 b1 + sin25 b2 = ( 0.114774,  0.917074, 0.381818),
%   d2 = -sin25 b1 + cos25 b2 = (-0.502129, -0.278094, 0.818855),
% where b1 = (1,3,0)/sqrt10 and b2 = (-9,3,20)/sqrt490 span v^perp, at half
% widths 2.00 along d1 and 1.05 along d2.  The plane shadow (I-P)e_3 has
% coordinates (0.381823, 0.818857) in that frame, so it lies inside the patch.
\begin{tikzpicture}[
  x=1.85cm, y=1.85cm,
  vec/.style   ={-{Latex[length=2mm,width=1.5mm]}, line width=0.75pt},
  proj/.style  ={line width=1.05pt, black!85},
  edge/.style  ={line width=0.45pt, black!65},
  hidden/.style={line width=0.45pt, black!45, dashed, dash pattern=on 1.4pt off 1.2pt},
  frame/.style ={densely dotted, line width=0.3pt, black!35},
  constr/.style={densely dotted, line width=0.45pt, black!55},
  hair/.style  ={line width=0.4pt, black!45},
  every node/.style={inner sep=1.5pt, font=\small}]

% ---- the kernel plane, as a rectangle through the origin -----------------
\coordinate (Cpp) at ( 1.3009, 0.9659);   %  2.00 d1 + 1.05 d2
\coordinate (Cpm) at ( 0.9683,-1.1025);   %  2.00 d1 - 1.05 d2
\coordinate (Cmm) at (-1.3009,-0.9659);   % -2.00 d1 - 1.05 d2
\coordinate (Cmp) at (-0.9683, 1.1025);   % -2.00 d1 + 1.05 d2
\path[fill=black!12] (Cpp)--(Cmp)--(Cmm)--(Cpm)--cycle;
\draw[edge]          (Cpp)--(Cmp)--(Cmm)--(Cpm)--cycle;

% ---- the three coordinate axes, as a frame of reference ------------------
\coordinate (O)  at (0,0);
\coordinate (A1) at (-0.7425,-0.3712);    % 1.05 e_1
\coordinate (A2) at ( 0.7425,-0.3712);    % 1.05 e_2
\coordinate (A3) at ( 0,      1.3423);    % 1.55 e_3
\draw[frame] (O)--(A1);
\draw[frame] (O)--(A2);
\draw[frame] (O)--(A3);
\node[anchor=north,      font=\footnotesize] at (A1) {$e_1$};
\node[anchor=north,      font=\footnotesize] at (A2) {$e_2$};
\node[anchor=south west, font=\footnotesize] at (A3) {$e_3$};

% ---- the range of the chart: the line R v --------------------------------
\coordinate (V)    at (-0.8081, 0.1691);  % v      = ( 6,-2, 3)/7
\coordinate (Vfar) at (-1.1718, 0.2452);  % 1.45 v
\coordinate (Vneg) at ( 0.3637,-0.0761);  % -0.45 v, behind the plane
\draw[hidden] (O)--(Vneg);
\draw[edge]   (V)--(Vfar);
% the line meets the plane at a right angle: the square on -0.17v and -0.17d2
\draw[edge, line width=0.4pt]
      ( 0.1374,-0.0288)--( 0.1105,-0.1962)--(-0.0269,-0.1674);

% ---- the two shadows of e_3, and the parallelogram they close ------------
\coordinate (Pe) at (-0.3463, 0.0725);    % P e_3    = ( 18,-6, 9)/49
\coordinate (Qe) at ( 0.3463, 0.7935);    % (I-P)e_3 = (-18, 6,40)/49
\coordinate (E3) at ( 0,      0.8660);    % e_3
\draw[constr] (Pe)--(E3) (Qe)--(E3);
\draw[proj]   (O)--(Pe);
\draw[proj]   (O)--(Qe);
\fill (Pe) circle (1.3pt);
\fill (Qe) circle (1.3pt);

% ---- the unit vector on the line, and the selected basis vector ----------
\draw[vec] (O)--(V);
\draw[vec] (O)--(E3);
\fill (O) circle (1.1pt);

\node[anchor=east,       font=\footnotesize] at (-1.20, 0.26) {$\operatorname{ran}P$};
\node[anchor=south,      font=\footnotesize] at (-0.68, 0.19) {$v$};
\node[anchor=east,       font=\scriptsize]   at (-0.50,-0.10) {$Pe_3$};
\node[anchor=north west, font=\scriptsize]   at ( 0.36, 0.60) {$(I-P)e_3$};
\node[font=\footnotesize] at (-0.30,-0.62) {$\ker P$};

% ======================= the forced normalization =========================
% Nine rows 0.32 apart, the two columns starting at 1.55 and 3.98.
\begin{scope}[every node/.style={anchor=west, inner sep=1.5pt,
                                 font=\footnotesize}]
\node[font=\scriptsize] at (1.55, 1.63)
  {$m=3$, \ $k=1$, \ $t=\tfrac1{49}(9,36,4)$, \ $g=(2,-\tfrac13,\tfrac32)$,
   \ $v=\tfrac17(6,-2,3)$};
\node[font=\scriptsize] at (1.55, 1.31)
  {$\lVert Pe_3\rVert^2=s_3=\tfrac9{49}$, \
   $\lVert (I-P)e_3\rVert^2=\lVert z_3\rVert^2=\tfrac{40}{49}$};
\node[font=\scriptsize] at (1.55, 0.99)
  {$C=\{3\}$, \ $E:=C^{\complement}=\{1,2\}$, \
   $W_{33}=s_3-t_3=\tfrac5{49}>0$: \ $C$ dominates};

\draw[hair] (1.55, 0.80)--(6.48, 0.80);
\draw[hair] (3.81, 0.77)--(3.81,-1.13);

\node at (1.55, 0.61) {co-weights (Theorem~\ref{thm:duality})};
\node at (3.98, 0.61) {weights (Remark~\ref{rem:naive})};
\node at (1.55, 0.29) {$w_c:=z_c/\sqrt{1-t_c}$};
\node at (3.98, 0.29) {$u_c:=z_c/\sqrt{t_c}$};
\node at (1.55,-0.03) {$h_c:=R\adj w_c$, \ $R\adj\Sdual R=I_2$};
\node at (3.98,-0.03) {chart $(ZZ\adj,t)=(I-P,t)$};
\node at (1.55,-0.35)
  {$\textstyle\sum_{c\in C}w_c^{\vphantom{\mathsf{H}}}w_c\adj\preceq I_2$};
\node at (3.98,-0.35)
  {$(I-P-D)[E]=\tfrac1{49}\begin{psmallmatrix}4&12\\12&9\end{psmallmatrix}$};
\node at (1.55,-0.67) {$\lVert w_3\rVert^2=\tfrac{1-s_3}{1-t_3}=\tfrac{40}{45}\le1$};
\node at (3.98,-0.67) {$\det=\tfrac{36-144}{49^2}<0$};
\node at (1.55,-1.03) {$E$ dominates, matching $C$};
\node at (3.98,-1.03) {$E$ fails, though $C$ dominates};
\end{scope}

\end{tikzpicture}

  \caption{Duality at $m=3$, $k=1$ (Theorem~\ref{thm:duality}).  The vector
  $v_c=\sqrt{t_c}\,g_c$ has norm one, so the chart $P=vv\adj$ projects onto a
  line of the index space and $I-P$ onto the complementary plane, where the
  shadow of $e_3$ is the kernel row $z_3$ in coordinates.  Selecting $\{3\}$
  from $P$ and deleting it from $I-P$ are one operation, and no weight enters
  it.  The weights enter the normalization.  Rescaling the kernel rows by the
  co-weights and whitening by $\Sdual$ carries the dominating singleton
  $\{3\}$ to the dominating pair $\{1,2\}$; the rescaling by the weights of
  Remark~\ref{rem:naive}, whose chart is $(I-P,t)$ itself, refuses that pair.
  Example~\ref{ex:uneven} carries the failure in the other direction at
  $m=2$.}
  \label{fig:duality}
\end{figure}

\begin{example}\label{ex:uneven}
Take $m=2$, $k=1$, $t=(\tfrac15,\tfrac45)$ and $g=(2,\tfrac12)$; the
Parseval condition holds, $\tfrac15\cdot4+\tfrac45\cdot\tfrac14=1$.  Here
$v=\bigl(2/\sqrt5,\;1/\sqrt5\bigr)$, of norm one, so
$\ker P$ is spanned by the unit vector $z=(-1/\sqrt5,\,2/\sqrt5)$, giving
$\lvert z_1\rvert^2=\tfrac15$ and $\lvert z_2\rvert^2=\tfrac45$.  The
singleton $\{1\}$ dominates, $g_1^2=4\ge1$, and $\{2\}$ does not,
$g_2^2=\tfrac14<1$.

The correct dual: $w_1^2=\tfrac{1/5}{4/5}=\tfrac14$,
$w_2^2=\tfrac{4/5}{1/5}=4$, so
$\Sdual=\tfrac15\cdot\tfrac14+\tfrac45\cdot4=\tfrac{13}{4}$ and
$h_c^2=w_c^2/\Sdual$, that is $h_1^2=\tfrac1{13}$ and $h_2^2=\tfrac{16}{13}$.
Parseval holds, $\tfrac15\cdot\tfrac1{13}+\tfrac45\cdot\tfrac{16}{13}=1$,
and the correspondence is exact: $\{2\}$ dominates in $\dsn^{\perp}$
because $\tfrac{16}{13}\ge1$, matching $\{1\}$ upstairs, and $\{1\}$
fails downstairs because $\tfrac1{13}<1$, matching $\{2\}$ upstairs.

The naive dual of Remark~\ref{rem:naive}: $u_1^2=\tfrac{1/5}{1/5}=1$ and
$u_2^2=\tfrac{4/5}{4/5}=1$.  Both singletons dominate, so the
correspondence fails at $C=\{2\}$, which does not dominate upstairs while
its complement does downstairs.
\end{example}

At $(4,2)$ the two cells of Corollary~\ref{cor:dualcell} coincide and the
correspondence becomes an involution on the six pairs.

\begin{example}\label{ex:dualworked}
Take the design of Example~\ref{ex:chartworked}, at $(4,2)$ with
$t=\tfrac19(2,3,3,1)$ and the chart printed there.  Since $m-k=k=2$, the
dual $\dsn^{\perp}$ is again a design at $(4,2)$ and
$C\mapsto C^{\complement}$ sends pairs to pairs.

The kernel of $P$ is orthogonal to the two columns $(2,2,1,0)$ and
$(0,1,-2,2)$ of $3V$.  The integer vectors
\[
  n_1=(-5,4,2,0),\qquad n_2=(0,-2,4,5)
\]
annihilate both: $2(-5)+2(4)+1(2)=0$ and $1(4)+(-2)(2)=0$ for $n_1$,
$2(-2)+1(4)=0$ and $1(-2)+(-2)(4)+2(5)=0$ for $n_2$.  They also satisfy
$\lVert n_1\rVert^2=\lVert n_2\rVert^2=45$ and $\langle n_1,n_2\rangle=0$,
so $Z:=\tfrac1{3\sqrt5}(n_1\ n_2)$ has $Z\adj Z=I_2$ and
$ZZ\adj=I_4-P$.  Its rows are
\[
  z_0=\tfrac1{3\sqrt5}(-5,0),\quad
  z_1=\tfrac1{3\sqrt5}(4,-2),\quad
  z_2=\tfrac1{3\sqrt5}(2,4),\quad
  z_3=\tfrac1{3\sqrt5}(0,5),
\]
of squared lengths $\tfrac19(5,4,4,5)$, which is $1-s_c$ index by index.

Rescale by the co-weights $1-t=\tfrac19(7,6,6,8)$.  Then
$\lVert w_c\rVert^2=\bigl(\tfrac57,\tfrac23,\tfrac23,\tfrac58\bigr)$ and
the operator \eqref{eq:dualgram} is
\[
  \Sdual=\sum_c t_c\,w_c^{\vphantom{\mathsf{H}}}w_c\adj
    =\diag\bigl(\tfrac8{21},\tfrac7{24}\bigr),
  \qquad \det\Sdual=\tfrac19,
\]
positive definite as Lemma~\ref{lem:dualgram} requires.  Take
$R=\diag\bigl(\sqrt{21/8},\sqrt{24/7}\bigr)$, so that $R\adj\Sdual R=I_2$
as in Lemma~\ref{lem:whiten}.  The dual atoms $h_c=R\adj w_c$ are
\[
  h_0=\tfrac14(-\sqrt{30},0),\quad
  h_1=\tfrac1{\sqrt{35}}(7,-4),\quad
  h_2=\tfrac1{2\sqrt{35}}(7,16),\quad
  h_3=\tfrac17(0,\sqrt{105}),
\]
and the radicals cancel again in the rank-one terms:
\[
  h_0^{\vphantom{\mathsf{H}}}h_0\adj
    =\begin{pmatrix}\tfrac{15}8&0\\[2pt]0&0\end{pmatrix},
  \quad
  h_1^{\vphantom{\mathsf{H}}}h_1\adj
    =\begin{pmatrix}\tfrac75&-\tfrac45\\[2pt]
                    -\tfrac45&\tfrac{16}{35}\end{pmatrix},
  \quad
  h_2^{\vphantom{\mathsf{H}}}h_2\adj
    =\begin{pmatrix}\tfrac7{20}&\tfrac45\\[2pt]
                    \tfrac45&\tfrac{64}{35}\end{pmatrix},
  \quad
  h_3^{\vphantom{\mathsf{H}}}h_3\adj
    =\begin{pmatrix}0&0\\[2pt]0&\tfrac{15}7\end{pmatrix}.
\]
Their $t$-weighted sum is $I_2$: the first diagonal entries give
$\tfrac5{12}+\tfrac7{15}+\tfrac7{60}=1$, the two off-diagonal
contributions cancel, and the second diagonal entries give
$\tfrac{16}{105}+\tfrac{64}{105}+\tfrac{25}{105}=1$.  So $\dsn^{\perp}$ is
a design at $(4,2)$ carrying the weights of $\dsn$, with leverage scores
$t_c\lVert h_c\rVert^2=\tfrac1{84}(35,52,61,20)$ summing to $2$.

Figure~\ref{fig:dualworked} runs the correspondence at all six pairs.  Take
two of them.  Upstairs $\{0,2\}$ dominates with slack
zero; downstairs its complement gives
\[
  h_1^{\vphantom{\mathsf{H}}}h_1\adj+h_3^{\vphantom{\mathsf{H}}}h_3\adj-I_2
   =\begin{pmatrix}\tfrac25&-\tfrac45\\[3pt]
                   -\tfrac45&\tfrac85\end{pmatrix},
\]
of determinant $0$ and trace $2$, so $\{1,3\}$ dominates in $\dsn^{\perp}$
with slack zero as well.  The spectra are $\{1,\tfrac83\}$ upstairs and
$\{1,3\}$ downstairs, and $\lmin=1$ at both ends.  Upstairs $\{2,3\}$
fails; downstairs
$h_0^{\vphantom{\mathsf{H}}}h_0\adj+h_1^{\vphantom{\mathsf{H}}}h_1\adj-I_2$
has $(2,2)$ entry $\tfrac{16}{35}-1<0$, so $\{0,1\}$ fails in
$\dsn^{\perp}$.

The naive dual of Remark~\ref{rem:naive} is also a design at $(4,2)$ with
the weights of $\dsn$, of chart $(I_4-P,t)$ and gap
\[
  (I_4-P)-D=\tfrac19\begin{pmatrix}
    3&-4&-2&0\\ -4&1&0&-2\\ -2&0&1&4\\ 0&-2&4&4\end{pmatrix} .
\]
Over $\{0,1\},\{0,2\},\{0,3\},\{1,2\},\{1,3\},\{2,3\}$ its six principal
$2\times2$ minors, multiplied by $81$, are $-13$, $-1$, $12$, $1$, $0$,
$-12$.  Five carry the sign the correspondence demands.  The sixth is
$E=\{0,2\}$, whose block
$\tfrac19\begin{psmallmatrix}3&-2\\-2&1\end{psmallmatrix}$ has
determinant $-\tfrac1{81}$, so $E$ fails there, while the complementary
pair $\{1,3\}$ dominates strictly in $\dsn$ with $81\det W[\{1,3\}]=2$.
\end{example}

\begin{figure}[ht]
  \centering
  % Duality at (m,k) = (4,2), the self-dual cell, on the design of
% Example ex:chartworked.  Three panels: the kernel rows z_c, the dual atoms
% h_c against the unit circle, and the two normalizations run against each
% other over all six pairs.  The cell is self-dual, m - k = k = 2, so the dual
% design lives at (4,2) as well and the complement map C -> C^c is an
% involution on the six pairs: the opposite-edge map of the square in
% figures/chartworked.tex.
%
% PROJECTION.  Both drawn panels are plane figures and carry no projection of a
% higher-dimensional space; the map is the identity of R^2, scaled by 1.55cm in
% the kernel panel and by 1.05cm in the dual-atom panel,
%   panel A:  X = 1.55 x_1,  Y = 1.55 x_2
%   panel B:  X = 1.05 x_1,  Y = 1.05 x_2
% so no viewing normal and no visibility rule are needed and no line is hidden.
% The two scales differ because the longest kernel row has norm
% sqrt(5/9) = 0.7453560 while the longest dual atom has norm
% sqrt(61/28) = 1.4759985, a ratio of sqrt(549/140) = 1.9802597; at a common
% scale one panel would be twice the other.  The unit circle is drawn in both,
% so the eye compares against it and not across panels.
%
% THE PRIMAL DATA.  Over R, m = 4, k = 2, t = (2,3,3,1)/9,
%   g_0 = (sqrt2,0), g_1 = (2,1)/sqrt3, g_2 = (1,-2)/sqrt3, g_3 = (0,2),
% with s = (4,5,5,4)/9 and the rational chart 9P = [[4,4,2,0],[4,5,0,2],
% [2,0,5,-4],[0,2,-4,4]].  See figures/chartworked.tex for the Parseval check.
%
% THE KERNEL BASIS.  ker P is orthogonal to the two columns (2,2,1,0) and
% (0,1,-2,2) of 3V.  The integer vectors n1 = (-5,4,2,0) and n2 = (0,-2,4,5)
% annihilate both -- 2(-5)+2(4)+1(2) = 0, 1(4)+(-2)(2) = 0, 2(-2)+1(4) = 0,
% 1(-2)+(-2)(4)+2(5) = 0 -- and satisfy n1.n1 = n2.n2 = 45, n1.n2 = 0.  So
% Z = [n1 n2]/(3 sqrt5) has Z^T Z = I_2 and Z Z^T = I - P, and its rows are
%   z_0 = (-5, 0)/(3 sqrt5)   |z_0|^2 = 5/9 = 1 - s_0
%   z_1 = ( 4,-2)/(3 sqrt5)   |z_1|^2 = 4/9 = 1 - s_1
%   z_2 = ( 2, 4)/(3 sqrt5)   |z_2|^2 = 4/9 = 1 - s_2
%   z_3 = ( 0, 5)/(3 sqrt5)   |z_3|^2 = 5/9 = 1 - s_3
% Panel A draws these four rows.
%
% THE TWO NORMALIZATIONS.  With 1 - t = (7,6,6,8)/9,
%   w_c = z_c/sqrt(1-t_c):  |w_c|^2 = (5/7, 2/3, 2/3, 5/8),
%   S^perp = sum_c t_c w_c w_c^T = diag(8/21, 7/24),  det = 1/9,
%   R = diag(sqrt(21/8), sqrt(24/7)),  R^T S^perp R = I_2,
%   h_0 = (-sqrt(30)/4, 0)          = (-sqrt30, 0)/4       |h_0|^2 = 15/8
%   h_1 = ( sqrt(35)/5, -4 sqrt(35)/35) = (7,-4)/sqrt35     |h_1|^2 = 13/7
%   h_2 = ( sqrt(35)/10, 8 sqrt(35)/35) = (7,16)/(2 sqrt35) |h_2|^2 = 61/28
%   h_3 = ( 0, sqrt(105)/7)         = (0, sqrt105)/7        |h_3|^2 = 15/7
% and sum_c t_c |h_c|^2 = 5/12 + 13/21 + 61/84 + 5/21 = 2 = m - k.  The naive
% rescaling u_c = z_c/sqrt(t_c) gives |u_c|^2 = (5/2, 4/3, 4/3, 5) and chart
% (Z Z^T, t) = (I - P, t), whose gap is 9(I-P-D) = [[3,-4,-2,0],[-4,1,0,-2],
% [-2,0,1,4],[0,-2,4,4]].  Panel B draws the h_c; the u_c are not drawn, since
% u_c is a positive multiple of z_c and panel A already carries their
% directions, all four of which the naive dual keeps.
%
% ELLIPSE FRAME.  At E = {1,3} the dual selection matrix is
% h_1 h_1^T + h_3 h_3^T = [[7/5,-4/5],[-4/5,13/5]], of spectrum {1,3}, so the
% ellipse {x : x^T (h_1 h_1^T + h_3 h_3^T) x = 1} has semi-axes 1 and
% 1/sqrt3 = 0.5773503, the long one along the eigenvector (2,1)/sqrt5 at
% atan2(1,2) = 26.5651 degrees, and it touches the unit circle exactly at
% +-(2,1)/sqrt5 = +-(0.8944272, 0.4472136).  Upstairs the same picture stands
% at C = {0,2} with spectrum {1, 8/3} and tangency along (1,2)/sqrt5.
%
% THE TABLE.  81 det W[C] over {0,1}, {0,2}, {0,3}, {1,2}, {1,3}, {2,3} is
% -12, 0, +6, +4, +2, -10; det(sum_{c in E} h_c h_c^T - I_2) over the
% complementary pairs {2,3}, {1,3}, {1,2}, {0,3}, {0,2}, {0,1} is -18/7, 0,
% 27/28, 1, 3/8, -15/8; and 81 det (I-P-D)[E] over the same is -12, 0, +1, +12,
% -1, -13.  The sign agrees in all six rows for the co-weight normalization and
% in five for the weight normalization, the exception being C = {1,3} against
% E = {0,2}, boxed in the body.
%
% DISCLOSED: this design is new to the paper and no Lean declaration mentions
% it; every number here was computed in exact arithmetic outside the assistant
% -- rational for the chart, S^perp, the leverage scores and all three minor
% columns, algebraic for z_c, R, h_c and the ellipse frame, whose decimals in
% the body are truncations of those exact values.  Theorem thm:duality is
% mechanized over R, but as an EXISTENTIAL over dual designs, and by a
% construction that is not the one drawn: the formal proof whitens the
% co-Parseval operator and takes the dual atoms from an orthonormal
% completion.  What panel B draws is therefore the printed
% construction, and the formal statement covers it only in asserting that some
% dual design with these weights flips domination.  No complex Naimark duality
% exists in the development, so nothing here transfers to C.
\begin{tikzpicture}[
  x=1cm, y=1cm,
  vec/.style   ={-{Latex[length=2mm,width=1.5mm]}, line width=0.75pt},
  rim/.style   ={line width=0.6pt, black!70},
  lev/.style   ={line width=0.5pt, black!70},
  frame/.style ={densely dotted, line width=0.3pt, black!35},
  tight/.style ={line width=0.5pt, black!55, dashed,
                 dash pattern=on 2.2pt off 1.6pt},
  hair/.style  ={line width=0.4pt, black!45},
  box/.style   ={line width=0.45pt, black!60},
  tag/.style   ={font=\footnotesize, fill=white, inner sep=0.8pt},
  every node/.style={inner sep=1.5pt, font=\small}]

% ==================== panel A: the rows of the kernel basis ================
\begin{scope}[shift={(2.00,3.90)}]
  \draw[frame] (-1.72,0)--(1.72,0);
  \draw[frame] (0,-1.72)--(0,1.72);
  \draw[rim] (0,0) circle (1.55);

  \draw[vec] (0,0)--(-1.15530, 0.00000);   % z_0 = (-5, 0)/(3 sqrt5)
  \draw[vec] (0,0)--( 0.92424,-0.46212);   % z_1 = ( 4,-2)/(3 sqrt5)
  \draw[vec] (0,0)--( 0.46212, 0.92424);   % z_2 = ( 2, 4)/(3 sqrt5)
  \draw[vec] (0,0)--( 0.00000, 1.15530);   % z_3 = ( 0, 5)/(3 sqrt5)
  \fill (0,0) circle (1.1pt);
  \node[tag, anchor=east]       at (-1.19, 0.00) {$z_0$};
  \node[tag, anchor=north west] at ( 0.95,-0.48) {$z_1$};
  \node[tag, anchor=south west] at ( 0.50, 0.92) {$z_2$};
  \node[tag, anchor=east]       at (-0.06, 1.10) {$z_3$};

  \node[anchor=south, font=\footnotesize] at (0,1.64)
    {$\ker P$, \ $Z\adj Z=I_2$};
  \node[anchor=north, align=center, font=\footnotesize] at (0,-1.64)
    {$\lVert z_c\rVert^2=1-s_c$\\[2pt] $=\tfrac19(5,4,4,5)$};
\end{scope}

% ==================== panel B: the dual atoms h_c ==========================
\begin{scope}[shift={(6.30,3.90)}]
  \draw[frame] (-1.72,0)--(1.72,0);
  \draw[frame] (0,-1.30)--(0,1.72);
  \draw[rim] (0,0) circle (1.05);
  % the ellipse of h_1 h_1^H + h_3 h_3^H, spectrum {1,3}: semi-axes 1 and
  % 1/sqrt3, the long one along (2,1)/sqrt5
  \draw[lev, rotate=26.5651] (0,0) ellipse (1.05 and 0.60622);
  \draw[tight] (-1.16454,-0.58227)--(1.16454,0.58227);       % +-1.24 (2,1)/sqrt5
  \draw[lev, fill=white] ( 0.93915, 0.46957) circle (1.4pt); % + (2,1)/sqrt5
  \draw[lev, fill=white] (-0.93915,-0.46957) circle (1.4pt); % - (2,1)/sqrt5

  \draw[vec] (0,0)--(-1.43777, 0.00000);   % h_0 = (-sqrt30/4, 0)
  \draw[vec] (0,0)--( 1.24238,-0.70993);   % h_1 = (sqrt35/5, -4 sqrt35/35)
  \draw[vec] (0,0)--( 0.62119, 1.41986);   % h_2 = (sqrt35/10, 8 sqrt35/35)
  \draw[vec] (0,0)--( 0.00000, 1.53704);   % h_3 = (0, sqrt105/7)
  \fill (0,0) circle (1.1pt);
  \node[tag, anchor=south]      at (-1.40, 0.05) {$h_0$};
  \node[tag, anchor=north west] at ( 1.26,-0.72) {$h_1$};
  \node[tag, anchor=west]       at ( 0.68, 1.38) {$h_2$};
  \node[tag, anchor=east]       at (-0.08, 1.38) {$h_3$};

  \node[anchor=south, font=\footnotesize] at (0,1.72)
    {$\dsn^{\perp}$, \
     $\textstyle\sum_ct_ch_c^{\vphantom{\mathsf{H}}}h_c\adj=I_2$};
  \node[anchor=north, align=center, font=\footnotesize] at (0,-1.34)
    {$x\adj\bigl(h_1^{\vphantom{\mathsf{H}}}h_1\adj
       +h_3^{\vphantom{\mathsf{H}}}h_3\adj\bigr)x=1$\\[2pt]
     spectrum $\{1,3\}$, tangent at $\pm\tfrac1{\sqrt5}(2,1)$};
\end{scope}

% =============== the two normalizations, side by side ======================
\begin{scope}[every node/.style={anchor=west, inner sep=1.5pt,
                                 font=\footnotesize}]
  \node at ( 8.90, 5.36) {co-weights};
  \node at (13.10, 5.36) {weights};
  % the header rule spans the block it underlines: from the west edge of the
  % left column's nodes at 8.90 to the east edge of the widest ink in the right
  % column, the 79.25pt of the u_c-norm row at 13.10 plus its 1.5pt inner sep,
  % which is 15.938; its node box reaches 15.991, so the rule adds no width.
  \draw[hair] ( 8.90, 5.18)--(15.94, 5.18);
  \draw[hair] (12.94, 5.16)--(12.94, 3.16);
  \node at ( 8.90, 4.96) {$w_c=z_c/\sqrt{1-t_c}$};
  \node at (13.10, 4.96) {$u_c=z_c/\sqrt{t_c}$};
  \node at ( 8.90, 4.56) {$\Sdual=\diag\bigl(\tfrac8{21},\tfrac7{24}\bigr)$};
  \node at (13.10, 4.56) {chart $(I-P,\,t)$};
  \node at ( 8.90, 4.16) {$R=\diag\bigl(\sqrt{21/8},\sqrt{24/7}\bigr)$};
  \node at (13.10, 4.10)
    {$\lVert u_c\rVert^2=\bigl(\tfrac52,\tfrac43,\tfrac43,5\bigr)$};
  \node at ( 8.90, 3.76) {$h_c=R\adj w_c$};
  \node at ( 8.90, 3.36)
    {$\lVert h_c\rVert^2=\bigl(\tfrac{15}8,\tfrac{13}7,
       \tfrac{61}{28},\tfrac{15}7\bigr)$};
\end{scope}

% ==================== the six pairs, three columns =========================
\begin{scope}[every node/.style={inner sep=1.5pt, font=\scriptsize}]
  \node[anchor=west] at ( 1.30, 0.72) {$C$};
  \node[anchor=east] at ( 4.10, 0.72) {$81\det W[C]$};
  \node[anchor=west] at ( 5.20, 0.72) {$E=C^{\complement}$};
  \node[anchor=east] at ( 9.60, 0.72)
    {$\det\bigl(\sum_{c\in E}h_c^{\vphantom{\mathsf{H}}}h_c\adj-I_2\bigr)$};
  \node[anchor=east] at (13.90, 0.72) {$81\det(I-P-D)[E]$};
  \draw[hair] ( 1.20, 0.44)--(14.00, 0.44);
  \draw[hair] ( 4.65, 0.92)--( 4.65,-2.30);
  \draw[hair] (10.20, 0.92)--(10.20,-2.30);

  % Rows 0.44 apart, in the order {0,1} {0,2} {0,3} {1,2} {1,3} {2,3}.  The
  % pitch is set by column four: a \tfrac at \scriptsize has height 7.137pt and
  % depth 3.137pt, so a pitch of 10.274pt = 0.3610 would put the denominator of
  % one fraction row on the numerator of the next, which 0.36 did; 0.44 =
  % 12.520pt leaves 2.25pt, the gap the reading block of
  % figures/chartworked.tex runs at.  The header sits at 0.72 rather than 0.66
  % so that its tallest entry, whose centred box reaches 6.014pt below the row,
  % clears the rule at 0.44 by 1.94pt instead of 0.24pt.
  \foreach \yy/\CC/\dW/\EE/\dH/\dN in {%
     0.14/{\{0,1\}}/{-12}/{\{2,3\}}/{-\tfrac{18}7}/{-12},
    -0.30/{\{0,2\}}/{0}/{\{1,3\}}/{0}/{0},
    -0.74/{\{0,3\}}/{+6}/{\{1,2\}}/{\tfrac{27}{28}}/{+1},
    -1.18/{\{1,2\}}/{+4}/{\{0,3\}}/{1}/{+12},
    -1.62/{\{1,3\}}/{+2}/{\{0,2\}}/{\tfrac38}/{-1},
    -2.06/{\{2,3\}}/{-10}/{\{0,1\}}/{-\tfrac{15}8}/{-13}}{%
    \node[anchor=west] at ( 1.30,\yy) {$\CC$};
    \node[anchor=east] at ( 4.10,\yy) {$\dW$};
    \node[anchor=west] at ( 5.20,\yy) {$\EE$};
    \node[anchor=east] at ( 9.60,\yy) {$\dH$};
    \node[anchor=east] at (13.90,\yy) {$\dN$};}
  % The one row where the weight normalization disagrees in sign.  The entry
  % $-1$ is 10.861pt wide with 1.5pt of inner sep, so its ink runs from 13.4658
  % to 13.8473 and, centred on the row, from 0.1052 above to 0.1052 below it;
  % the frame clears that by 1.87, 2.07 and 2.13pt on the left, the right and
  % the two horizontals.  The previous frame cleared the left by 0.26pt.
  \draw[box] (13.40,-1.80) rectangle (13.92,-1.44);

  \node[anchor=west, font=\footnotesize] at ( 1.20,-2.60)
    {The sign of column two matches column four at all six pairs and
     column five at five.};
\end{scope}

\end{tikzpicture}

  \caption{Duality at $(4,2)$, where the dual cell is the cell itself, on
  the design of Example~\ref{ex:chartworked}.  Left: the rows of an
  orthonormal basis of $\ker P$, of squared lengths
  $1-s_c=\tfrac19(5,4,4,5)$.  Centre: the dual atoms of
  Theorem~\ref{thm:duality}, and the ellipse of
  $h_1^{\vphantom{\mathsf{H}}}h_1\adj+h_3^{\vphantom{\mathsf{H}}}h_3\adj$,
  of spectrum $\{1,3\}$, touching the unit circle at
  $\pm\tfrac1{\sqrt5}(2,1)$.  Right: the two normalizations of
  Remark~\ref{rem:naive}.  Below: at each of the six pairs, the gap minor of
  $\dsn$ and the minor of each normalization, the boxed entry the one at which
  the rescaling by the weights takes the wrong sign.
  The dual atoms drawn are those of the construction above
  Theorem~\ref{thm:duality}; the mechanized statement asserts only that some
  dual design with these weights flips domination.}
  \label{fig:dualworked}
\end{figure}

\subsection{Corank one}\label{sec:corankone}

Let $m=k+1$ with $k\ge1$.  Then $\ker P$ is a line, spanned by a unit
vector $z\in\FF^{m}$, and $P=I-zz\adj$.  Set
\begin{equation}\label{eq:pvars}
  p_c\;:=\;\frac{\lvert z_c\rvert^2}{1-t_c}
  \;=\;\frac{1-t_c\ell_c}{1-t_c},
\end{equation}
the second expression by $\lvert z_c\rvert^2=1-P_{cc}=1-s_c$
(Proposition~\ref{prop:chartid}).  Since $z$ is a unit vector,
\begin{equation}\label{eq:pnorm}
  \sum_c p_c\,(1-t_c)\;=\;\sum_c\lvert z_c\rvert^2\;=\;1,
  \qquad\text{while}\qquad
  \sum_c (1-t_c)\;=\;m-1\;=\;k .
\end{equation}
The $p_c$ therefore have co-weighted mean $1/k$.  Domination compares a
positive diagonal with a rank-one perturbation of it, and one scalar
decides such a comparison.  Figure~\ref{fig:rankoneball} draws the
comparison.

\begin{lemma}[Lean]\label{lem:rankone}
Let $A\pd0$ and let $b$ be a vector, and put $q:=b\adj A^{-1}b$.  Then
\[
  A-bb\adj\psd0\iff q\le1,
  \qquad
  A-bb\adj\pd0\iff q<1 .
\]
\end{lemma}

\begin{proof}
Put $y:=A^{-1}b$, so that $Ay=b$ and
$q=\langle y,b\rangle=\langle y,Ay\rangle$: the number $q$ is the squared
length of $y$ in the inner product $\langle\cdot,A\cdot\rangle$, and in
particular $q\ge0$.  The form of the perturbed matrix is
$\langle x,(A-bb\adj)x\rangle=\langle x,Ax\rangle-\lvert\langle
b,x\rangle\rvert^2$.

Suppose $q\le1$.  The form $\langle\cdot,A\cdot\rangle$ is an inner
product, and $\langle b,x\rangle=\langle Ay,x\rangle=\langle y,Ax\rangle$
is the $A$-inner product of $y$ and $x$; Cauchy--Schwarz in it gives
$\lvert\langle b,x\rangle\rvert^2\le q\,\langle x,Ax\rangle$.  Hence
$\langle x,(A-bb\adj)x\rangle\ge(1-q)\langle x,Ax\rangle$, which is
nonnegative, and strictly positive for $x\ne0$ when $q<1$.

Conversely, evaluate at $x=y$: the form equals $q-q^2=q(1-q)$.  If
$A-bb\adj\psd0$ this is nonnegative, which forces $q\le1$ since $q\ge0$.
If moreover $q=1$ then $y\ne0$, because $\langle y,Ay\rangle=1$, and the
form vanishes at $y$, so $A-bb\adj$ is not positive definite.
\end{proof}

\begin{figure}[ht]
  \centering
  % The rank-one criterion as a point-in-ellipsoid test, drawn at three
% coordinates.  A = diag(4, 9/4, 1) is positive definite, so the locus q = 1 is
% the ellipsoid with semi-axes 2, 3/2, 1, and three vectors b run from the
% origin: (1, 3/4, 0) with q = 1/2, (1, 0, sqrt3/2) with q = 1, (0, 3/2, 1) with
% q = 2.  Each of the three lies in a coordinate plane, so the point where its
% ray meets the ellipsoid -- the open circle, at 1/sqrt(q) times b -- sits on a
% drawn principal ellipse.  Orthographic projection of R^3 with azimuth 20 and
% elevation 22 degrees, scaled by 2.20cm:
%   X = 2.20(-0.342020 x_1 + 0.939693 x_2)
%   Y = 2.20(-0.352015 x_1 - 0.128123 x_2 + 0.927184 x_3)
% Every coordinate below is the parenthesised value, with its source triple in a
% trailing comment; the 2.20cm is carried by the x and y unit vectors.  A curve
% is given by the projections of two conjugate semi-axis vectors, so that
% (X, Y) = (cos u) proj(first) + (sin u) proj(second).
\begin{tikzpicture}[
  x=2.20cm, y=2.20cm,
  vec/.style={-{Latex[length=2mm,width=1.5mm]}, line width=0.75pt},
  frame/.style={densely dotted, line width=0.3pt, black!35},
  edge/.style={line width=0.45pt, black!65},
  rim/.style={line width=0.45pt, black!70},
  hidden/.style={line width=0.45pt, black!45, dashed, dash pattern=on 1.4pt off 1.2pt},
  slack/.style={densely dotted, line width=0.5pt, black!55},
  hit/.style={line width=0.5pt, black!85},
  lead/.style={line width=0.3pt, black!55},
  every node/.style={font=\small}]

% ---- the three coordinate axes, as a frame of reference ------------------
% each runs from the origin past its pole, which sits on the ellipsoid
\draw[frame] (0,0) -- (-0.807167,-0.830755);   % (2.36, 0,     0    )
\draw[frame] (0,0) -- ( 1.620970,-0.221012);   % (0,    1.725, 0    )
\draw[frame] (0,0) -- ( 0,        1.038446);   % (0,    0,     1.12 )

% ---- the ellipsoid: three principal ellipses, hidden arcs first ----------
% A point x of the ellipsoid faces the viewer when (A^{-1}d).x > 0, with
% d = (0.871268, 0.317116, 0.374607) the viewing direction; that inequality
% cuts each principal ellipse at the two parameters recorded below.
% E_12, semi-axis vectors (2,0,0) -> (-0.684040,-0.704030)
%                     and (0,3/2,0) -> ( 1.409539,-0.192185)
\draw[hidden, variable=\u, domain=115.887:295.887, samples=61, smooth]
  plot ({-0.684040*cos(\u)+1.409539*sin(\u)}, {-0.704030*cos(\u)-0.192185*sin(\u)});
% E_13, semi-axis vectors (2,0,0) and (0,0,1) -> (0, 0.927184)
\draw[hidden, variable=\u, domain=130.693:310.693, samples=61, smooth]
  plot ({-0.684040*cos(\u)}, {-0.704030*cos(\u)+0.927184*sin(\u)});
% E_23, semi-axis vectors (0,3/2,0) and (0,0,1)
\draw[hidden, variable=\u, domain=150.562:330.562, samples=61, smooth]
  plot ({1.409539*cos(\u)}, {-0.192185*cos(\u)+0.927184*sin(\u)});
\draw[edge, variable=\u, domain=-64.113:115.887, samples=61, smooth]
  plot ({-0.684040*cos(\u)+1.409539*sin(\u)}, {-0.704030*cos(\u)-0.192185*sin(\u)});
\draw[edge, variable=\u, domain=-49.307:130.693, samples=61, smooth]
  plot ({-0.684040*cos(\u)}, {-0.704030*cos(\u)+0.927184*sin(\u)});
\draw[edge, variable=\u, domain=-29.438:150.562, samples=61, smooth]
  plot ({1.409539*cos(\u)}, {-0.192185*cos(\u)+0.927184*sin(\u)});
% the silhouette: the contour generator is the section of the ellipsoid by the
% plane (A^{-1}d).x = 0, and its projection is the ellipse through
% (1.566752, 0.134478) and (0, 1.172254) at the two conjugate parameters,
% equivalently the ellipse of semi-axes 1.579547 and 1.162758 turned by
% 10.8171 degrees
\draw[rim, rotate=10.8171] (0,0) ellipse (1.579547 and 1.162758);

% ---- the three vectors, each with the point where its ray leaves ----------
% q = 1/2: the ray leaves at sqrt2 times b, on E_12 at parameter 45 degrees
\draw[vec]   (0,0) -- (0.362749,-0.448107);                 % b = (1, 3/4, 0)
\draw[slack] (0.362749,-0.448107) -- (0.513005,-0.633719);  % (1.414214, 1.060660, 0)
\draw[hit]   (0.513005,-0.633719) circle (1.7pt);
% q = 1: the tip is itself the exit point, on E_13 at parameter 60 degrees
\draw[vec] (0,0) -- (-0.342020,0.450950);                   % b = (1, 0, sqrt3/2)
\draw[hit] (-0.342020,0.450950) circle (1.7pt);
% q = 2: the shaft leaves at b over sqrt2, on E_23 at parameter 45 degrees
\draw[vec] (0,0) -- (1.409539,0.734999);                    % b = (0, 3/2, 1)
\draw[hit] (0.996695,0.519723) circle (1.7pt);              % (0, 1.060660, 0.707107)
\fill (0,0) circle (1.1pt);

% ---- labels ---------------------------------------------------------------
\node[anchor=north east, font=\footnotesize] at (-0.83,-0.85) {$b_1$};
\node[anchor=west,       font=\footnotesize] at ( 1.66,-0.23) {$b_2$};
\node[anchor=east,       font=\footnotesize] at (-0.04, 1.05) {$b_3$};
\draw[lead] (-0.60,1.140) -- (-0.536,1.056);   % to the rim at parameter 110 degrees
\node[anchor=east, align=right, font=\footnotesize] at (-0.62,1.19)
  {$b\adj A^{-1}b=1$};
\node[align=center, font=\footnotesize] at ( 0.80,-0.31) {$q=\tfrac12$\\[-1pt] inside};
\node[align=center, font=\footnotesize] at (-0.74, 0.36) {$q=1$\\[-1pt] on};
\node[anchor=west, align=left, font=\footnotesize] at ( 1.46, 0.80)
  {$q=2$\\[-1pt] outside};

% ---- the reading, below the picture ---------------------------------------
\node[anchor=north, align=left, font=\scriptsize] at (0.14,-1.50) {%
  $A=\diag(4,\tfrac94,1)$,\quad $q=b\adj A^{-1}b$,\quad
  semi-axes $\sqrt{A_{11}},\sqrt{A_{22}},\sqrt{A_{33}}=2,\tfrac32,1$\\[4pt]
  \begin{tabular}{@{}l@{\hspace{1.3em}}l@{\hspace{1.3em}}l@{\hspace{1.3em}}l@{}}
    $b=(1,\tfrac34,0)$          & $q=\tfrac12$ & $A-bb\adj\pd0$
      & $C$ dominates strictly\\[1pt]
    $b=(1,0,\tfrac{\sqrt3}{2})$ & $q=1$        & $A-bb\adj\psd0$, singular
      & $C$ dominates, not strictly\\[1pt]
    $b=(0,\tfrac32,1)$          & $q=2$        & $A-bb\adj$ indefinite
      & $C$ does not dominate\\
  \end{tabular}\\[4pt]
  at $m=k+1$, with $C=\{1,\dots,m\}\setminus\{c_0\}$: \
  $A=\diag(1-t_c)_{c\in C}$, \ $b=(z_c)_{c\in C}$, \ $q=\sum_{c\ne c_0}p_c$};

\end{tikzpicture}

  \caption{Lemma~\ref{lem:rankone} in three coordinates.  The locus $q=1$ is the
  image of the unit sphere under $A^{1/2}$, and the criterion asks which side of
  it $b$ lies on; the open circle marks where the ray through $b$ leaves the
  ellipsoid.  At $m=k+1$ the comparison is \eqref{eq:corankone}, the ratio at the
  omitted index against the $t$-weighted mean of all of them.  The ellipsoid then
  lives in the index space: its coordinates are the selected indices, its
  semi-axes are the square roots of the co-weights $1-t_c$, and the vectors $g_c$
  do not enter.  Figure~\ref{fig:jensen} is the same comparison in bar-chart
  form.}
  \label{fig:rankoneball}
\end{figure}

\begin{theorem}[Lean]\label{thm:corankone}
Let $m=k+1$, $k\ge1$, over either field.
\begin{enumerate}
\item[(i)] The complement of $\{c_0\}$ dominates if and only if
\begin{equation}\label{eq:corankone}
  p_{c_0}\;\ge\;\sum_{c} t_c\,p_c ,
\end{equation}
and it dominates strictly if and only if the inequality is strict.
\item[(ii)] $\Gtz(k+1,k)$ and $\GtzC(k+1,k)$ hold.
\item[(iii)] Equality holds in \eqref{eq:corankone} for every $c_0$ if
and only if $p_c\equiv1/k$, if and only if
\begin{equation}\label{eq:tielaw}
  k\,t_c\,\ell_c \;=\; k-1+t_c \qquad (1\le c\le m),
\end{equation}
and in that case every $k$-subset dominates and none dominates strictly.
\item[(iv)] For every weight vector $t$ on $k+1$ atoms there is a design
with weights exactly $t$ satisfying \eqref{eq:tielaw}.
\end{enumerate}
\end{theorem}

\begin{proof}
(i) Put $C:=\{1,\dots,m\}\setminus\{c_0\}$, of size $k$.  By
Lemma~\ref{lem:chart} the subset $C$ dominates precisely when
$W[C]=(I-zz\adj-D)[C]\psd0$, that is when $A-bb\adj\psd0$ with
$A:=\diag(1-t_c)_{c\in C}$ and $b:=(z_c)_{c\in C}$.  Every $t_c<1$, so
$A\pd0$ and $A^{-1}=\diag\bigl((1-t_c)^{-1}\bigr)_{c\in C}$, whence
\[
  q\;=\;b\adj A^{-1}b\;=\;\sum_{c\ne c_0}\frac{\lvert
  z_c\rvert^2}{1-t_c}\;=\;\sum_{c\ne c_0}p_c\;=\;\sum_c p_c-p_{c_0}.
\]
Lemma~\ref{lem:rankone} turns domination into $\sum_cp_c-p_{c_0}\le1$,
and strict domination into the strict inequality.  Expanding
\eqref{eq:pnorm} as $\sum_cp_c-\sum_ct_cp_c=1$ and substituting yields
\eqref{eq:corankone}.

(ii) The numbers $t_c$ are nonnegative and sum to one, so the weighted
mean $\sum_ct_cp_c$ does not exceed $\max_cp_c$.  Any index $c_0$
attaining the maximum satisfies \eqref{eq:corankone}, and its complement
is a dominating $k$-subset.  No step used the field.

(iii) Equality at every $c_0$ says that every $p_{c_0}$ equals the fixed
number $\sum_ct_cp_c$, so $p$ is constant, say $p_c\equiv\bar p$; then
\eqref{eq:pnorm} gives $\bar pk=1$.  Conversely a constant $p$ makes both
sides of \eqref{eq:corankone} equal.  For the second equivalence, by
\eqref{eq:pvars},
\[
  p_c=\frac1k
  \iff k\,(1-t_c\ell_c)=1-t_c
  \iff k\,t_c\ell_c=k-1+t_c ,
\]
which is \eqref{eq:tielaw}; equivalently
$\ell_c=(k-1+t_c)/(k\,t_c)$.  At corank one every $k$-subset is the
complement of a singleton, so by (i) equality everywhere means that every
$k$-subset dominates and none does so strictly.

(iv) Given $t$, the system \eqref{eq:tielaw} prescribes
$\lvert z_c\rvert^2=(1-t_c)/k$, and these numbers sum to $k/k=1$ by
\eqref{eq:pnorm}; so put
\[
  z_c\;:=\;\sqrt{\frac{1-t_c}{k}}\;>\;0 .
\]
It remains to realize $P=I-zz\adj$ by a design with weight vector $t$,
which is Lemma~\ref{lem:realize}; explicitly, let $e_1$ be the first
standard basis vector and $u:=z-e_1$, nonzero because
$z_1^2=(1-t_1)/k<1$, and put
\[
  Q\;:=\;I-\frac{2\,uu\adj}{\lVert u\rVert^2},
  \qquad \lVert u\rVert^2=2(1-z_1),
\]
the Householder reflection.  It is Hermitian and involutive, hence
unitary, and $Qe_1=e_1+u=z$ because $u\adj e_1=z_1-1$.  Its first column
is thus $z$, so its remaining $k$ columns are an orthonormal basis of
$z^{\perp}$; let $V$ be the $m\times k$ matrix formed by them and
$g_c:=t_c^{-1/2}v_c$ as in Lemma~\ref{lem:realize}.  Then $V\adj V=I_k$,
and expanding $I=QQ\adj$ as the sum of the outer products of the columns
of $Q$ gives $zz\adj+VV\adj=I$.  So the design so obtained has weight
vector $t$ and chart $P=VV\adj=I-zz\adj$, whose diagonal gives
$t_c\ell_c=1-(1-t_c)/k$, that is \eqref{eq:tielaw}.
\end{proof}

No step of the argument sees the field, so the complex collapse of
\S\ref{sec:complex} must wait one size longer.  Part (iv) produces a
family of designs attaining the threshold with no slack, one over each
point of the open weight simplex; \S\ref{sec:ties} draws the consequence.
Figure~\ref{fig:jensen} shows the two cases.

\begin{example}\label{ex:tetrachart}
The tetrahedron of Example~\ref{ex:tetra} is the uniform member at $k=3$.
There $t_c=\tfrac14$ and \eqref{eq:tielaw} reads
$3\cdot\tfrac14\cdot\ell_c=3-1+\tfrac14=\tfrac94$, so $\ell_c=3$, as
computed there.  The kernel vector is
$z=\bigl(\tfrac12,\tfrac12,\tfrac12,\tfrac12\bigr)$, since
$\lvert z_c\rvert^2=(1-\tfrac14)/3=\tfrac14$, and the chart is
\[
  P\;=\;I_4-\tfrac14J_4,
  \qquad
  P_{cc}=\tfrac34=s_c,
  \qquad
  P_{cd}=-\tfrac14=\tfrac14\langle g_c,g_d\rangle\ (c\ne d),
\]
where $J_4$ is the all-ones matrix; $P^2=I-\tfrac12J+\tfrac1{16}J^2
=I-\tfrac12J+\tfrac14J=P$ and $\tr P=3$.  Every $p_c$ equals
$\tfrac{1/4}{3/4}=\tfrac13=1/k$, so \eqref{eq:corankone} holds with
equality at every $c_0$; by part (iii) every triple dominates and none
dominates strictly, which is the spectrum $\{1,4,4\}$ recorded in
Example~\ref{ex:tetra}.
\end{example}

\begin{remark}\label{rem:ktwo}
At $k=2$ the law \eqref{eq:tielaw} reads
$\ell_c=\tfrac12+\tfrac1{2t_c}$.  Its uniform-weight shadow is classical:
for $t_c=p/n$ the right-hand side is $\tfrac12+\tfrac{n}{2p}$, and these
are the magnitudes appearing in the equality description of
\cite[Proposition~1]{Nesterenko2026real}, whose three clusters of sizes
$p+q+r=n$ correspond to the three vectors of a corank-one design with
weights $p/n$, $q/n$, $r/n$.
\end{remark}

\begin{figure}[ht]
  \centering
  % The corank-one mechanism at m = k+1.  Left: the pivots p_c of a design whose
% pivot vector is not constant; the maximal pivot lies above the t-weighted
% mean, and the complement of a maximising index dominates.  Right: the
% extremal case, in which every pivot equals 1/k.
\begin{tikzpicture}[
  x=1cm, y=1cm,
  bar/.style={draw, line width=0.5pt, fill=black!10},
  hit/.style={draw, line width=0.7pt,
              pattern=north east lines, pattern color=black!55},
  mean/.style={dashed, dash pattern=on 1.9pt off 1.6pt, line width=0.55pt},
  every node/.style={font=\small}]

% ================= p not constant =========================================
\begin{scope}
  % heights are 4.4 p_c for p = (1/2, 1/4, 2/5, 1/5); the weighted mean is 7/20
  \draw[hit] (0.20,0) rectangle (0.70,2.20);
  \draw[bar] (0.95,0) rectangle (1.45,1.10);
  \draw[bar] (1.70,0) rectangle (2.20,1.76);
  \draw[bar] (2.45,0) rectangle (2.95,0.88);
  \draw[line width=0.6pt] (0.02,0)--(3.13,0);

  \draw[mean] (0.05,1.54)--(3.28,1.54);
  \node[anchor=west, font=\footnotesize] at (3.32,1.54) {$\sum_c t_cp_c$};

  \node[font=\footnotesize] at (0.45,2.42) {$c_0$};
  \foreach \x/\lab in {0.45/1, 1.20/2, 1.95/3, 2.70/4}
    {\node[anchor=north, font=\footnotesize] at (\x,-0.06) {$p_{\lab}$};}

  \node at (1.575,2.78) {$p$ not constant};
  \node[font=\footnotesize, align=center] at (1.575,-1.02)
    {$p_{c_0}\ \ge\ \sum_c t_cp_c$\\[2pt]
     $\{1,\dots,m\}\setminus\{c_0\}$ dominates};
\end{scope}

% ================= p constant =============================================
\begin{scope}[shift={(6.05,0)}]
  % every pivot equals 1/k with k = 3
  \draw[mean] (0.05,1.4667)--(3.28,1.4667);
  \node[anchor=west, font=\footnotesize] at (3.32,1.4667) {$1/k$};
  \foreach \x in {0.20, 0.95, 1.70, 2.45}
    {\draw[bar] (\x,0) rectangle (\x+0.50,1.4667);}
  \draw[line width=0.6pt] (0.02,0)--(3.13,0);

  \foreach \x/\lab in {0.45/1, 1.20/2, 1.95/3, 2.70/4}
    {\node[anchor=north, font=\footnotesize] at (\x,-0.06) {$p_{\lab}$};}

  \node at (1.575,2.78) {$p$ constant};
  \node[font=\footnotesize, align=center] at (1.575,-1.02)
    {$k\,t_c\ell_c=k-1+t_c$\\[2pt]
     every $k$-subset dominates, none strictly};
\end{scope}

\end{tikzpicture}

  \caption{Theorem~\ref{thm:corankone} at $(4,3)$.  The ratios $p_c$ of
  \eqref{eq:pvars} have co-weighted mean $1/k$ by \eqref{eq:pnorm}.  Left:
  the largest $p_{c_0}$ is at least the $t$-weighted mean $\sum_ct_cp_c$,
  so by \eqref{eq:corankone} its complement dominates.  Right: when $p$ is
  constant every $p_c$ equals $1/k$, the tie law \eqref{eq:tielaw} holds,
  every $k$-subset dominates and none dominates strictly.}
  \label{fig:jensen}
\end{figure}

Corank two is one further step: Corollary~\ref{cor:dualcell} converts it
into rank two, which is Theorem~\ref{thm:ranktwo}, and \S\ref{sec:corank}
carries out the conversion once that input is proved.
